% Readout revision, 2026-10-02: based on the user's Downloads manuscript.
% Adds alpha/beta security theory to Results, Methods, Extended Data Fig. 7
% and a subsection within Supplementary Note 2.
\documentclass[fleqn,10pt]{wlscirep}
\usepackage[utf8]{inputenc}
\usepackage[T1]{fontenc}
\usepackage{bm}
\usepackage{graphicx}
\graphicspath{{figures/}}

\usepackage{booktabs}
\usepackage{adjustbox}
\usepackage{longtable}
\usepackage{caption}
\usepackage{mathtools}
\usepackage{amssymb}
\usepackage{amsmath}
\usepackage{etoolbox}
\usepackage{amsthm}
\usepackage{tocloft}
\usepackage{float}
\usepackage{titlesec}

\newtheorem{proposition}{Proposition}
\newtheorem{theorem}{Theorem}
\newtheorem{lemma}{Lemma}
\newtheorem{corollary}{Corollary}
\newtheorem{remark}{remark}
% \usepackage{algorithm}
% \usepackage{algorithmic}
% \usepackage[ruled,vlined,linesnumbered]{algorithm2e}
\DeclareMathAlphabet{\mathcal}{OMS}{cmsy}{m}{n}

\usepackage{algorithm}
\usepackage{algorithmicx}
\usepackage{algpseudocode}

\usepackage{lineno}
%\linenumbers

\allowdisplaybreaks

% \renewcommand{\qedsymbol}{} % 重新定义qedsymbol为空，这样末尾就不会显示方框



\title{Indirect reciprocity with inertial reputation updating}

\author[1,2]{Wenqiang Zhu}
\author[1,2*]{Other Authors}
\affil[1]{School of Artificial Intelligence, Beihang University, Beijing, 100191, China}
\affil[2]{Key laboratory of Mathematics, Informatics and Behavioral Semantics, Beihang University, Beijing 100191, China}
\affil[*]{Corresponding to: 11~(1.~1); 22 (2.~2)}


\begin{abstract}
Reputations support cooperation by linking past behavior to future treatment. How should they change when new evidence challenges an established judgment? We model inertial reputation updating through separate pathways for losing good standing and recovering it under private assessment. Across four second-order social norms, suitable configurations improve strategy discrimination, reduce disagreement and expand the basins leading to cooperation. Equal total memory depths produce markedly different outcomes depending on their allocation between the two pathways and the social norm. These gains can coexist with vulnerability to occasional under-cooperation. In representative configurations, clustering the same average exploitation into persistent episodes closes invasion windows. Allowing intermediate warning states to influence behavior can also prevent invasion, but is constrained by resident cooperation and the existence of a high-reputation state. Our results show how the direction of reputational persistence and the use of intermediate evidence jointly shape the capacity of social norms to sustain cooperation. 
\end{abstract}


\begin{document}

% \flushbottom
\raggedbottom
\maketitle

% \thispagestyle{empty}


\maketitle


\section*{Introduction}

% P1. Behavioural importance and the central question.
Reputations enable past behavior to influence future cooperation. By helping others, individuals can become more likely to receive help from third parties, sustaining indirect reciprocity even without repeated encounters between the same partners~\cite{nowak1998,wedekind2000,schmid2021}. This mechanism depends on how observers revise their judgments as new behavior becomes visible~\cite{michelmata2024}. A new observation may reinforce an established reputation, but it may also contradict it. When new evidence conflicts with an existing reputation, how should that reputation change? Understanding this process is necessary to explain how reputations shape the evolution of cooperation in human societies.

% P2. An illustrative social situation, followed by bounded empirical support.
Consider a colleague known for helping others who declines a request on one occasion. An observer may register the lapse while continuing to regard that colleague as helpful. Conversely, a single helpful act by someone previously seen as unhelpful may prompt reconsideration without immediately restoring a favorable reputation. Experiments show that prior good deeds can temper judgments of subsequent misconduct, depending on its ambiguity and context~\cite{effron2010}, and that moral beliefs can be revised differently for people previously judged good or bad~\cite{siegel2018}. Research on trust repair also shows that reliable behavior over repeated interactions can restore damaged trust, with recovery affected by prior deception~\cite{schweitzer2006}. These findings motivate a process in which new evidence can modify a judgment without immediately reversing it, and in which retaining a good reputation and recovering from a bad one need not operate identically.

% P3. The binary baseline.

Widely studied models represent reputation as a binary label, good or bad, and use a social norm to update that label after an observed interaction~\cite{nowak2005,ohtsuki2004,hilbe2018}. Second-order norms base the assessment on the donor's action and the recipient's reputation ~\cite{pacheco2006,ohtsuki2007}; third-order norms also use the donor's current reputation~\cite{brandt2004,ohtsuki2006}. The latter therefore allow past assessments to influence a new judgment. Yet if the current good-or-bad label is the entire reputational record, it cannot indicate whether that judgment has already survived conflicting evidence. If an unfavorable observation leaves someone labeled good, another observation under the same circumstances is processed from the same good state, with the same probability of changing the label. Likewise, a bad label does not reveal whether favorable evidence has already begun to accumulate. To examine how reputations respond to repeated challenges, we need to track evidence that can change before the good-or-bad label does.


% P4. Private assessment and the closest existing mechanisms.
Even individuals who share the same social norm can disagree about whether a particular act deserves approval when reputations are private. A donor may refuse help to someone they regard as bad, while an observer who regards the recipient as good condemns that refusal. These judgments shape subsequent help and refusal, which become evidence for further judgments; disagreement about one person can therefore affect the reputations of others~\cite{hilbe2018,fujimoto2022,kessinger2023}. Research has examined several ways to limit the consequences of these differences in perspective. Public institutions can provide a common source of reputational judgments, while gossip allows private opinions to spread and become more aligned~\cite{radzvilavicius2021,kawakatsu2024}. Observers can also take the donor's perspective into account when interpreting an action~\cite{radzvilavicius2019}, or use the donor's treatment of a recipient as evidence for revising the recipient's reputation, as in dual private assessment~\cite{tham2026}. A complementary approach changes how each observer forms a judgment from the evidence available. Cumulative image scores and additional reputation states have long provided ways to retain information beyond a single binary assessment~\cite{nowak1998,tanabe2013,murase2022}. Recent work has shown how this richer record can sustain cooperation under private information. Quantitative assessment updates a score after each evaluation and converts it into a good or bad label, allowing errors to be corrected through subsequent experience~\cite{schmid2023}. Aggregating several observations with some tolerance for negative assessments can likewise generate agreement and support cooperation without requiring observers to share their judgments~\cite{michelmata2024}. These advances show why an established reputation need not be overturned by one contrary observation. They also bring the subsequent process of revision into focus. After an apparent lapse, later conduct may reaffirm the earlier judgment or challenge it repeatedly. Preserving a good reputation can protect cooperation from a mistaken assessment, but may also delay recognition of sustained uncooperative behavior. Preserving a bad reputation can require further evidence before help is restored, but may also prolong the exclusion of a wrongly condemned cooperator. The consequences of persistence therefore depend on how new evidence strengthens or reverses an existing judgment. Understanding when such persistence supports cooperation requires examining both how good reputations are lost and how bad reputations are repaired, and how these processes interact under a given social norm.

% P5. A concrete model response, including the distinction from a binary
Here we develop a framework of inertial reputation updating to examine how the persistence of good and bad judgments shapes indirect reciprocity under private assessment. Observers retain internal states that record evidence challenging an established reputation. Consecutive negative assessments move a good record through warning states towards bad standing, while consecutive positive assessments move a bad record along a recovery path towards good standing. An assessment consistent with the established reputation resets this accumulation. The two paths can require different amounts of evidence, allowing tolerance of negative assessments and strictness about rehabilitation to vary separately under the same social norm. We also allow intermediate states to be read as good or bad with different probabilities when guiding behavior and moral assessment. New evidence can therefore influence social interactions before either path is completed. 

Across four standard second-order social norms, we identify configurations that improve the distinction between cooperative strategies and unconditional defectors, reduce private disagreement and enlarge the basins leading to cooperative behavior. The effective allocation of memory depends on the norm, and configurations with the same total depth can produce substantially different cooperative outcomes. We further show that improved recognition of unconditional defectors can coexist with vulnerability to occasional under-cooperation. The resulting invasion opportunities depend on the temporal organization of exploitation and on how intermediate evidence is used. Together, these results connect the persistence of individual judgments to the incentives supporting cooperation, showing that the cooperative consequences of a social norm depend on how reputations are preserved and revised.

\section*{Results}\label{sec_model}

\subsection*{A framework for inertial reputation updating in indirect reciprocity}
\label{sec:model-overview}

\begin{figure}[!ht]
\centering
\includegraphics[width=0.98\linewidth]{Fig1_model.pdf}
\caption{\textbf{A framework for inertial reputation updating.}
\textbf{a}, A donor chooses whether to cooperate (C) or defect (D), and an observer assesses the realized action as positive ($+$) or negative ($-$) according to a social norm.
\textbf{b}, A DISC donor acts on its own image of the recipient, whereas the observer assesses the action using its independently held image of that recipient. These private images can disagree.
\textbf{c}, Intended cooperation becomes defection with probability $e_1$, and the prescribed assessment is reversed with probability $e_2$.
\textbf{d}, Instantaneous updating replaces the donor's image with the latest assessment.
\textbf{e}, One-step tolerance inertia introduces a warning state ($\mathrm W$) between established good ($\mathrm G$) and bad ($\mathrm B$) records. Two consecutive negative assessments change $\mathrm G$ to $\mathrm B$; a positive assessment clears the warning. State $\mathrm W$ is read as good with probability $\alpha$.
\textbf{f}, One-step strictness inertia introduces state $\mathrm S$ along the recovery path. Two consecutive positive assessments change $\mathrm B$ to $\mathrm G$; a negative assessment resets this path. State $\mathrm S$ is read as good with probability $1-\beta$. Binary readouts guide both behavior and subsequent moral assessment.}
\label{fig:model-overview}
\end{figure}

We consider a population whose members repeatedly interact in a donation game, a standard model of indirect reciprocity~\cite{nowak2005}. A donor can cooperate, paying a cost $c$ to provide a benefit $b>c$ to a recipient, or defect, providing no benefit and paying no cost (Fig.~\ref{fig:model-overview}a). Unconditional cooperators (ALLC) intend to help every recipient, unconditional defectors (ALLD) never intend to help, and discriminators (DISC) intend to help recipients they regard as good. Reputations are private: a donor and an observer may hold different views of the same recipient, so an action that the donor considers justified may receive a negative assessment from the observer~\cite{hilbe2018} (Fig.~\ref{fig:model-overview}b). Intended cooperation fails with probability $e_1$, and an assessment is recorded with the opposite sign with probability $e_2$ (Fig.~\ref{fig:model-overview}c). We examine Scoring, Shunning, Stern Judging and Simple Standing, which differ in how they assess actions towards recipients regarded as bad (Supplementary Table~\ref{tab:si-norms}).

Under instantaneous updating (without inertia), the latest assessment replaces the observer's binary image of the donor (Fig.~\ref{fig:model-overview}d). We introduce inertia by retaining an internal record of evidence that challenges an established reputation. In the one-step model, a negative assessment of a good individual first produces a warning state, $\mathrm W$. A second negative assessment changes the record to bad, whereas a positive assessment restores it to good (Fig.~\ref{fig:model-overview}e). Conversely, a positive assessment of a bad individual produces an intermediate state, $\mathrm S$, along the recovery path. A second positive assessment restores good standing, whereas a negative assessment returns the record to bad (Fig.~\ref{fig:model-overview}f).

These internal states determine the binary image consulted during both action choice and moral assessment. A warning is read as good with probability $\alpha$, while a recovery state is read as good with probability $1-\beta$. Larger $\alpha$ preserves a good image after negative evidence; larger $\beta$ delays the restoration of a good image after positive evidence. Full inertia, $\alpha=\beta=1$, preserves the established image until the corresponding path is completed. Setting $\alpha=\beta=0$ recovers instantaneous binary updating. The multi-step extension introduces $k_1$ warning states and $k_2$ recovery states, with good-readout probabilities $\alpha^m$ and $1-\beta^n$ at steps $m$ and $n$, respectively. Evidence supporting the established judgment resets the entire corresponding path. Thus, path depth controls the evidence required for reversal, while readout controls how intermediate evidence affects current interactions.

We first examine reputation formation at fixed strategy composition, then use stationary reputations to calculate payoffs and the evolution of strategy frequencies. This evolutionary analysis assumes that reputations adjust faster than strategies; the updating process and stationary calculation are specified in Methods.


\subsection*{Inertia supports high reputation in homogeneous discriminator populations}
\label{sec:results-all-disc-reputation}

\begin{figure}[!ht]
\centering
\includegraphics[width=\linewidth]{Fig2_reputation dynamics.pdf}
\caption{\textbf{One-step inertia supports a high-reputation branch under four social norms.}
\textbf{a--d}, Reputation equilibria in homogeneous DISC populations under Scoring (\textbf{a}), Shunning (\textbf{b}), Stern Judging (\textbf{c}) and Simple Standing (\textbf{d}). Curves show the right-hand side of the auxiliary reputation-equilibration equation in Methods, Eq.~\eqref{eq:methods-reputation-dynamics}, against $r_{\mathrm{DISC}}$. Grey denotes instantaneous updating, $(\alpha,\beta)=(0,0)$, and blue denotes full one-step inertia, $(\alpha,\beta)=(1,1)$. Zeros are self-consistent reputation states; filled and open circles classify their stability in this equilibration equation, and arrows show its direction. Full-state branch checks are described in Supplementary Note~\ref{sec:si-full-state-dynamics}. Parameters are $k_1=k_2=1$ and $e_1=e_2=0.02$.}
\label{fig:all-disc-reputation-rates}
\end{figure}

We first asked whether conditional cooperators could maintain good reputations when every individual used DISC. In this homogeneous population, the probability $r_{\mathrm{DISC}}$ of being regarded as good also determines intended cooperation. A stationary reputation must be consistent with the assessments generated by the resulting interactions. We calculated these states using the reputation mapping and self-consistency equations in Methods, Eqs.~\eqref{eq:methods-reputation-map} and~\eqref{eq:methods-reputation-dynamics}.

Full one-step inertia supported a high-reputation equilibrium under all four norms, with $r_{\mathrm{DISC}}^*\geq0.995$ at $e_1=e_2=0.02$ (Fig.~\ref{fig:all-disc-reputation-rates}). This contrasts with instantaneous updating, for which equilibrium reputation ranged from approximately $0.020$ under Shunning to $0.833$ under Simple Standing. On the high-reputation branch, donors and observers usually regard recipients as good, so DISC cooperation usually receives a positive assessment. An isolated negative assessment moves the record into a warning state without immediately changing its good readout, and a subsequent positive assessment clears that warning. This protection allows a high level of good reputation to remain self-consistent despite errors.

The high-reputation branch coexisted with lower-reputation states under Scoring, Shunning and Stern Judging. Simple Standing remained monostable at the displayed parameters (Fig.~\ref{fig:all-disc-reputation-rates}). The one-step $(\alpha,\beta)$ scan showed that bistability occupied a broad region under Scoring and arose at relatively high tolerance under Shunning and Stern Judging; no bistability was found under Simple Standing in this scan (Extended Data Fig.~\ref{fig:ed-bistability-alpha-beta}). Thus, the appearance of a high-reputation branch does not imply that every initial reputation state reaches it. Subsequent evolutionary comparisons use reputations selected from initially good records, with branch checks described in Methods and Supplementary Note~\ref{sec:si-full-state-dynamics}. These homogeneous-population results establish that high reputation can be maintained; we next ask whether reputations remain informative when behavioral strategies differ.


\subsection*{Directional memory improves recognition in mixed populations}
\label{sec:results-recognition}

\begin{figure}[!ht]
\centering
\includegraphics[width=\linewidth]{Fig3_reputation recognition and disagreement.pdf}
\caption{\textbf{The allocation of inertia shapes private agreement and strategy recognition.}
Columns show Scoring, Shunning, Stern Judging and Simple Standing.
\textbf{a--d}, Illustrative expressed-reputation matrices with $k_1=k_2=1$. Rows correspond to $(\alpha,\beta)=(0,0),(1,0),(0,1)$ and $(1,1)$. Within each row, matrices show DISC observers' images of DISC, ALLC and ALLD targets, from left to right. Blue denotes good and white bad.
\textbf{e--h}, Depth scans at $\alpha=\beta=1$. Cell color shows reputational separation $\Delta g$ and circle area shows mean disagreement $\bar d$ (Methods). Black and red outlines mark the binary baseline $(0,0)$ and full one-step inertia $(1,1)$. Stars mark the smallest total depths satisfying $\Delta g\geq0.8$ and $\bar d\leq0.1$, with ties resolved by greater separation and then lower disagreement. The selected pairs are $(2,0),(3,0),(4,3)$ and $(1,2)$ in column order; arrows connect the one-step reference to these pairs. Quantitative values are means of 30 independent runs, whereas panels \textbf{a--d} are separately saved illustrative snapshots. All simulations use $N=150$, with 50 individuals per strategy, initially good internal records and $e_1=e_2=0.02$.}
\label{fig:finite-population-reputation-dynamics}
\end{figure}

\begin{figure}[!ht]
\centering
\includegraphics[width=\linewidth]{Fig4_strategy dynamics.pdf}
\caption{\textbf{Inertia changes cooperative attraction basins and resistance to invasion.}
\textbf{a--d}, Strategy dynamics under Scoring, Shunning, Stern Judging and Simple Standing. Rows show the binary baseline, full one-step inertia, and representative full multi-step inertia at $(k_1,k_2)=(2,0),(3,0),(4,3)$ and $(1,2)$, respectively. Arrows indicate strategy change after reputation equilibration. Shading marks initial compositions reaching intended cooperation of at least $0.5$; percentages give intended cooperation at the highlighted equilibria. Filled and open circles denote stable and unstable equilibria; for an edge equilibrium, the classification refers to stability within that invariant edge.
\textbf{e--h}, Local invasion outcomes of the selected pure-DISC resident across implementation and assessment errors for the corresponding multi-step configurations. Rose denotes invasion by ALLD, purple denotes exclusion of ALLD but invasion by ALLC, and green denotes exclusion of both. Dashed and solid white boundaries mark zero growth of ALLD and ALLC, respectively. All panels use $b/c=5$ and initially good reputation records; panels \textbf{a--d} use $e_1=e_2=0.02$, and panels \textbf{e--h} use $\alpha=\beta=1$.}
\label{fig:strategy-dynamics}
\end{figure}

High reputation among discriminators is useful for cooperation only if defectors do not receive similarly favorable treatment. We therefore simulated private reputation formation in a finite population containing equal numbers of DISC, ALLC and ALLD, with strategy frequencies held fixed. One-step inertia had contrasting effects across norms (Fig.~\ref{fig:finite-population-reputation-dynamics}a--d). Under Scoring, full inertia preserved the good reputation of ALLC and the bad reputation of ALLD, but left considerable variation in the reputation of DISC. Under Shunning, nearly everyone was regarded as bad. Stern Judging retained substantial disagreement and little separation between strategies. Simple Standing generated a clearer distinction between cooperative strategies and ALLD. Agreement and recognition therefore need to be examined together: near-unanimous bad judgments, as under Shunning, can convey little information about behavioral differences.

We quantified these properties using reputational separation, $\Delta g$, between cooperative strategies and ALLD, and mean observer disagreement, $\bar d$, about the same targets (Methods, Eqs.~\eqref{eq:methods-recognition-separation} and~\eqref{eq:methods-recognition-disagreement}). With full readout, $\alpha=\beta=1$, varying the two path depths produced regions with both strong separation and low disagreement under every norm (Fig.~\ref{fig:finite-population-reputation-dynamics}e--h; Extended Data Fig.~\ref{fig:ed-finite-population-reputation-dynamics}). The effective allocation differed across norms. Scoring and Shunning benefited mainly from longer warning paths. Stern Judging benefited from retaining evidence along both paths. Under Simple Standing, extending the recovery path reduced the favorable reputations occasionally obtained by ALLD while preserving those of cooperative strategies.

To carry compact examples into the evolutionary analysis, we selected the shortest configurations in the surveyed grid satisfying $\Delta g\geq0.8$ and $\bar d\leq0.1$. The resulting pairs were $(k_1,k_2)=(2,0)$ for Scoring, $(3,0)$ for Shunning, $(4,3)$ for Stern Judging and $(1,2)$ for Simple Standing. Across 30 independent runs, these configurations produced mean separation between $0.881$ and $0.919$ and mean disagreement between $0.056$ and $0.081$. Full one-step inertia did not meet both criteria under any norm, although Simple Standing was close, with $\Delta g=0.815$ and $\bar d=0.108$. Changing the working thresholds changed some selected pairs (Extended Data Fig.~\ref{fig:ed-selection-robustness}). These configurations therefore provide norm-specific representatives for the next analysis; their selection is based on reputation formation, independently of their evolutionary performance.


\subsection*{Improved recognition expands the routes to cooperation}
\label{sec:results-strategy-dynamics}

Reliable recognition does not by itself establish whether cooperation is favored by selection. We therefore allowed ALLC, ALLD and DISC frequencies to evolve according to replicator dynamics, using payoffs evaluated after reputations had equilibrated (Methods). Full one-step inertia supported a highly cooperative DISC equilibrium under all four norms, but its basin of attraction remained limited under Shunning and Stern Judging. The representative multi-step configurations selected from the recognition analysis enlarged the cooperative basins and supported highlighted equilibria with at least $99.99\%$ intended cooperation at the reference parameters (Fig.~\ref{fig:strategy-dynamics}a--d). These percentages describe intended actions and implementation errors reduce the frequency of realized cooperation by the factor $1-e_1$.

The invasion conditions of an established DISC population explain another aspect of these outcomes. Let $r_{\mathrm R}^{*}$ be the resident good-reputation probability, and let $r_{\mathrm D}^{*}=\mathcal R\!\left[r_{\mathrm R}^{*}P^+_{\mathrm D\mid\mathrm G}+(1-r_{\mathrm R}^{*})P^+_{\mathrm D\mid\mathrm B}\right]$ and $r_{\mathrm C}^{*}=\mathcal R\!\left[r_{\mathrm R}^{*}P^+_{\mathrm C\mid\mathrm G}+(1-r_{\mathrm R}^{*})P^+_{\mathrm C\mid\mathrm B}\right]$ be the corresponding reputations of rare ALLD and ALLC mutants in that resident environment. When both reputation differences in the denominators are positive, pure DISC resists both mutants if
\begin{equation}
\frac{r_{\mathrm R}^{*}}{r_{\mathrm R}^{*}-\mathcal R\!\left[r_{\mathrm R}^{*}P^+_{\mathrm D\mid\mathrm G}+(1-r_{\mathrm R}^{*})P^+_{\mathrm D\mid\mathrm B}\right]}
<\frac bc<
\frac{1-r_{\mathrm R}^{*}}{\mathcal R\!\left[r_{\mathrm R}^{*}P^+_{\mathrm C\mid\mathrm G}+(1-r_{\mathrm R}^{*})P^+_{\mathrm C\mid\mathrm B}\right]-r_{\mathrm R}^{*}}.
\label{eq:results-disc-viability-window}
\end{equation}
The lower bound requires the loss of help experienced by ALLD to outweigh its saved cooperation cost. The upper bound reflects the converse possibility that ALLC gains enough additional help to offset its greater expenditure on cooperation. Boundary cases and the conditions when a finite bound is absent are given in Methods. These are local conditions on the selected DISC resident, distinct from the range of initial compositions that reach a cooperative state.

The parameter scans show how inertia changes these two outcomes. In the one-step model, increasing tolerance generally expanded cooperative basins, whereas strictness alone did not restore cooperation from interior populations under Scoring, Shunning or Stern Judging (Extended Data Fig.~\ref{fig:ed-strategy-alpha-beta}). The representative multi-step configurations also supported resistance to ALLD across broad regions of the error plane; under Scoring, Shunning and Stern Judging, substantial regions additionally excluded ALLC (Fig.~\ref{fig:strategy-dynamics}e--h). Their resistance to both classical mutants extended across much wider benefit-to-cost intervals than for full one-step inertia (Extended Data Fig.~\ref{fig:ed-strategy-bc}). Under Simple Standing, high cooperation could instead persist in a DISC--ALLC mixture after pure DISC lost resistance to ALLC, a possibility also present without inertia. Thus, inertia can improve the accessibility of cooperation and its resistance to invasion, with the relationship between these effects depending on the norm.


\subsection*{The cooperative value of memory depends on its allocation}
\label{sec:results-strategy-depth}

To separate the amount of memory from its direction, we compared all allocations with the same total depth, $K=k_1+k_2$. We summarized their evolutionary outcomes by averaging terminal cooperation over a common set of initial strategy compositions (Methods). At a fixed $K$, different allocations produced widely different cooperation levels (Fig.~\ref{fig:strategy-depth}). Increasing depth could improve cooperation when added to one path and reduce it when added to the other. Total depth alone therefore did not determine the evolutionary value of inertia.

The allocation patterns depended on the norm. Scoring and Shunning achieved their largest observed cooperation values with tolerance-only paths. Stern Judging benefited from retaining evidence in both directions, while Simple Standing favored configurations with a longer recovery path and a shorter warning path. These patterns reflect the different evidence generated by the norms. Scoring assigns negative assessments to refusals even when DISC regards the recipient as bad, making protection against reputational loss particularly useful. Under Simple Standing, ALLD can receive positive assessments by refusing recipients whom observers regard as bad; requiring further evidence before rehabilitation can reduce these favorable reputations. The full depth scans confirm that cooperation depends jointly on the two paths and that high cooperation and local protection of pure DISC occupy only partly overlapping regions (Extended Data Fig.~\ref{fig:ed-strategy-depth-all}). The configurations identified from recognition also yielded high cooperative returns at their respective depths, linking the informational and evolutionary consequences of how evidence is retained.

\begin{figure}[!ht]
\centering
\includegraphics[width=\linewidth]{Fig5_inertia allocation.pdf}
\caption{\textbf{Equal total memory depths can produce different cooperative outcomes.}
\textbf{a--d}, Basin-weighted realized cooperation against total depth $K=k_1+k_2$ under Scoring, Shunning, Stern Judging and Simple Standing. Small circles show the allocations at each $K$; color denotes the tolerance fraction $k_1/K$. Large outlined circles and their labels identify the largest observed terminal cooperation estimate at each depth, joined by the black line. The binary baseline at $K=0$ has no directional allocation. Each value averages over 512 common initial strategy compositions sampled uniformly from the simplex, with reputations selected from initially good records. Parameters are $b/c=5$, $e_1=e_2=0.02$ and $\alpha=\beta=1$. The stored terminal estimates for $(1,1)$ under Scoring, Shunning and Simple Standing did not pass the observable-convergence check; these points, including the Simple Standing envelope point at $K=2$, remain provisional.}
\label{fig:strategy-depth}
\end{figure}

\subsection*{Resistance to exploitation depends on behavioral timing and reputation readout}
\label{sec:results-opportunism}

\begin{figure}[!ht]
\centering
\includegraphics[width=\linewidth]{Fig6_strategy invasion.pdf}
\caption{\textbf{Partial exploitation depends on behavioral persistence and warning-state readout.}
Columns show Scoring, Shunning, Stern Judging and Simple Standing.
\textbf{a--d}, Normalized invasion exponent $\Lambda$ of an under-cooperation mutant that independently omits intended DISC cooperation with probability $\delta$. Grey, orange, blue and green curves show the binary baseline, full one-step inertia, representative full multi-step inertia, and the same representative depths with $\alpha=0.8$, respectively. The representative pairs are $(2,0),(3,0),(4,3)$ and $(1,2)$ in column order. Pale curves show additional depth configurations. Insets resolve small exploitation probabilities, and open circles mark non-zero endpoints of invasion intervals.
\textbf{e--h}, Invasion exponent for a mutant switching between DISC and ALLD modes, with stationary exploitative fraction $\delta$ and persistence $\eta$ between successive effective observations. These panels use the representative depths with full readout, $\alpha=\beta=1$; $\eta=0$ recovers the blue curves above. Warm colors indicate $\Lambda>0$, blue indicates $\Lambda<0$, and black contours mark $\Lambda=0$. The boundary $\eta=1$ represents the limit $\eta\to1^{-}$. All calculations use the pure-DISC resident selected from initially good records, $b/c=5$ and $e_1=e_2=0.02$. Graded readout uses $\beta=1$. The temporal observation probability is $p_{\mathrm{obs}}=1/149$; its relation to persistence in donor-action time is given in Methods, Eq.~\eqref{eq:methods-observed-persistence}. The plotted quantity is $\Lambda=(\Pi_{\mathrm X}-\Pi_{\mathrm R})/[(1-e_1)c]$.}
\label{fig:strategy-invasion}
\end{figure}

The preceding comparisons test cooperation against strategies that always cooperate or always defect. We next introduced an \emph{under-cooperation mutant} that follows DISC but independently withholds each intended cooperation with probability $\delta$. This parameter controls the long-run mean exploitation probability, with $\delta=0$ recovering DISC and $\delta=1$ recovering ALLD. In a pure-DISC resident population, let $r_{\mathrm X}^{*}(\delta)$ be the mutant's stationary good-reputation probability. Its normalized invasion exponent is
\begin{equation}
\Lambda(\delta)
=\frac{\Pi_{\mathrm X}-\Pi_{\mathrm R}}{(1-e_1)c}
=\delta r_{\mathrm R}^{*}
-\frac bc\left[r_{\mathrm R}^{*}-r_{\mathrm X}^{*}(\delta)\right].
\label{eq:results-partial-invasion}
\end{equation}
The mutant gains by saving cooperation costs and loses when its reduced reputation limits help from residents. Positive $\Lambda$ permits invasion. Resistance to ALLD tests only the endpoint $\delta=1$; exclusion of the under-cooperation family requires non-positive growth throughout $0<\delta\leq1$.

Full one-step inertia resisted ALLD under all four norms but left narrow invasion windows near $\delta=0$ (Fig.~\ref{fig:strategy-invasion}a--d). The representative multi-step configurations widened these windows under Scoring, Shunning and Stern Judging, to approximately $0<\delta<0.426$, $0.578$ and $0.365$, respectively. The Simple Standing configuration excluded the tested under-cooperation family. Thus, configurations that improved recognition of unconditional defectors and expanded cooperative basins could still leave weak behavioral deviations insufficiently penalized.

This vulnerability can be understood from the response of reputation near the resident state. Let $\Phi_{\mathrm R}^{*}$ be the resident positive-assessment probability and $\Delta\Phi^{*}$ the difference between positive-assessment probabilities for intended cooperation and defection in that environment. Independent omissions give $r_{\mathrm X}^{*}(\delta)=\mathcal R(\Phi_{\mathrm R}^{*}-\delta r_{\mathrm R}^{*}\Delta\Phi^{*})$, so
\begin{equation}
\Lambda'(0)
=r_{\mathrm R}^{*}
\left[1-\frac bc\,\Delta\Phi^{*}
\mathcal R'(\Phi_{\mathrm R}^{*})\right].
\label{eq:results-local-gain}
\end{equation}
A weak marginal reputation response allows a small reduction in cooperation to save more than it costs in lost help. Buffering isolated negative assessments can therefore protect both cooperative residents and occasional deliberate omissions. This local condition explains invasion close to the resident strategy; the complete windows follow from the full nonlinear reputation response.

We then held $\delta$ fixed and changed the temporal arrangement of exploitation. A \emph{temporally persistent opportunist} alternates between DISC and ALLD modes according to a two-state Markov process, spending a fraction $\delta$ of time in the exploitative mode. Observed persistence $\eta$ measures the correlation between modes at successive effective observations; $\eta=0$ recovers independent under-cooperation (Methods). In the representative full-readout configurations, increasing persistence contracted the invasion regions under Scoring, Shunning and Stern Judging and closed them at approximately $\eta=0.437$, $0.589$ and $0.527$ (Fig.~\ref{fig:strategy-invasion}e--g). Simple Standing remained resistant throughout the displayed range. Because average saved cost is unchanged, these differences arise from reputation. Persistent episodes are more likely to deliver successive negative assessments, while dispersed omissions are often interrupted by assessments that reset warnings. The calculation retains correlations between behavioral mode and internal reputation state; averaging assessment probabilities across modes would generally lose this effect.

Graded readout restored sensitivity to the first signs of under-cooperation. With full readout, an isolated negative assessment changes the internal record but leaves the individual treated as good. When $\alpha<1$, the first warning reduces the probability of receiving help before the established reputation is overturned. In an error-free, all-good DISC resident population, with $\beta=1$, at least one warning state and any recovery depth, the entire independent under-cooperation family has non-positive invasion fitness if and only if
\begin{equation}
\alpha\leq1-\frac cb.
\label{eq:results-alpha-security-benchmark}
\end{equation}
The first warning then produces a reputation response of magnitude $1-\alpha$, sufficient to offset the saved cooperation cost. This benchmark gives $\alpha=0.8$ a direct incentive-based interpretation at $b/c=5$; finite errors require evaluating the full stationary invasion criterion (Methods and Supplementary Note~\ref{sec:si-readout-security}). At the error rates used in Fig.~\ref{fig:strategy-invasion}, the largest secure readouts were $\alpha_{\mathrm{sec}}\simeq0.807$, $0.806$ and $0.804$ for the representative Scoring, Shunning and Stern Judging configurations. Simple Standing already excluded the tested family at $\alpha=1$. Setting $\alpha=0.8$ therefore closed all four under-cooperation windows while retaining resident intended cooperation between $0.978$ and $0.990$ (Extended Data Fig.~\ref{fig:ed-alpha-invasion}).

The same readout also changes the resident assessment environment, so it could not be reduced without affecting cooperation. Under Shunning and Stern Judging, the high-reputation stationary branch terminated at $\alpha_{\mathrm{rep}}\simeq0.591$ and $0.567$, respectively. The interval $\alpha_{\mathrm{rep}}<\alpha\leq\alpha_{\mathrm{sec}}$ retained that branch and excluded the tested mutants. These boundaries answer different questions: $\alpha_{\mathrm{rep}}$ follows from reputation self-consistency and is independent of $b/c$ at fixed norm, memory structure, $\beta$ and errors, whereas $\alpha_{\mathrm{sec}}$ depends on the payoff ratio. Under Scoring, cooperation instead declined continuously as $\alpha$ was reduced; Simple Standing had no corresponding branch fold. Thus, exclusion of under-cooperation does not by itself ensure high resident cooperation.

The recovery readout $\beta$ provides a distinct control. It lowers the probability that intermediate recovery states are treated as good, changing the reputation penalty after accumulated negative evidence. It has no effect in the representative Scoring and Shunning configurations because they contain no recovery states. In Stern Judging $(4,3)$, increasing $\beta$ narrowed the finite-exploitation window at $\alpha=1$ but did not remove invasion by arbitrarily weak under-cooperation. In Simple Standing $(1,2)$, sufficiently strong recovery strictness excluded the entire family even at $\alpha=1$. These configuration-specific effects follow from the different state occupancies and reputation responses, rather than a universal benefit of increasing $\beta$ (Supplementary Note~\ref{sec:si-readout-beta}).

The same invasion test also exposed partial exploitation of tolerance in representative Aggregation and Quantitative mechanisms~\cite{michelmata2024,schmid2023} (Extended Data Fig.~\ref{fig:ed-aggregation-quantitative}). Under Stern Judging, the tolerant Aggregation rule that accepts one positive assessment out of two sampled interactions admitted $0<\delta<0.137$. A Quantitative rule with scores from $-5$ to $5$ and a good-reputation threshold of zero admitted $0<\delta<0.383$. Both can preserve a good image despite some negative evidence, but the resulting protection of under-cooperation differs. Their temporal responses also differed: persistence closed the representative Quantitative window, whereas Aggregation was unchanged under the specified protocol of independently resampling a stationary infinite interaction pool. These comparisons concern the stated parameters and observation rules. Together, they show how the protection afforded by tolerant evidence processing can be exploited, and how its extent depends on the way assessments are retained and used.


\section*{Discussion}
How reputations change determines how past behavior influences
future opportunities for cooperation. Here we examined this process
by allowing evidence to accumulate along separate pathways for the
loss and recovery of reputation. Across four social norms, suitable
configurations reduced private disagreement, improved the distinction
between cooperative strategies and unconditional defectors, and
expanded the basins leading to cooperation. These effects depended
on how memory was allocated between the two pathways. Even with the
same total depth, different allocations produced substantially
different outcomes. The cooperative consequences of a social norm
therefore depend on the persistence and revision of the judgments
through which it operates.

The dependence on social norms follows from the feedback between
reputation, behavior and assessment. A reputation determines whom
a discriminator helps, while the observer's reputation of the
recipient helps determine whether that action is approved. Preserving
a judgment consequently changes both subsequent interactions and
the evidence generated by them. Under Scoring, for example, refusing
help receives a negative assessment even when the donor regards
the recipient as bad. Tolerance can protect discriminators from
the reputational consequences of these refusals. Under Simple
Standing, by contrast, defection towards a recipient regarded as
bad can receive a positive assessment, allowing unconditional
defectors to obtain favorable evidence. Requiring further positive
assessments before rehabilitation can then improve discrimination.
These differences help explain why the allocations associated with
high cooperation vary across norms. They also extend the study of
social norms by showing how a fixed assessment rule can produce
different social outcomes when embedded in different updating
processes \cite{ohtsuki2006,kessinger2023}.

Our results distinguish several requirements that are often discussed
together. High reputation in a homogeneous cooperative population
does not establish reliable recognition in a mixed population.
Agreement among observers can likewise coexist with poor recognition
if nearly everyone is regarded as bad. Moreover, favorable
reputational distinctions do not by themselves establish evolutionary
stability. We therefore examined reputation levels, disagreement,
strategy separation, attraction basins and invasion conditions
separately. The configurations identified from recognition criteria
also supported substantial cooperative outcomes in the evolutionary
analysis, linking these levels without making them equivalent.
Simple Standing illustrates the distinction particularly clearly:
a cooperative mixture of discriminators and unconditional cooperators
can persist even when pure discrimination is vulnerable to
unconditional cooperation.

This framework complements existing approaches to imperfect private
information. Empathetic assessment, public institutions, gossip and
dual assessment alter how individuals obtain, interpret or coordinate
social judgments
\cite{radzvilavicius2019,radzvilavicius2021,kawakatsu2024,tham2026}.
Aggregation and quantitative assessment show that retaining several
evaluations can support cooperation
\cite{michelmata2024,schmid2023}.
Our framework makes the directional persistence of established
judgments explicit and allows it to be varied while holding the
underlying norm fixed. Internal states distinguish individuals with
the same expressed reputation but different accumulated evidence
against that judgment. A fixed assessment rule using only current
binary labels cannot generally distinguish these records, although
the process is Markovian in its expanded state space. Some special
cases share stationary reputation mappings with existing aggregation
rules. The contribution lies in connecting the two directions of
reputational change to recognition, evolutionary outcomes and
responses to temporally structured behavior.

The invasion analysis reveals a further distinction between
recognizing different strategies and maintaining incentives against
small behavioral deviations. A system may assign very different
reputations to discriminators and unconditional defectors while
responding only weakly to a slight reduction in cooperation.
Under several representative configurations, the buffering that
preserved good reputations also allowed occasional under-cooperation
to save costs without a sufficient loss of help. The local invasion
condition identifies the relevant quantity as the marginal response
of stationary reputation to a change in assessment probability.
A large reputational difference between behavioral extremes need
not imply a strong response near the cooperative resident.
Accordingly, resistance to unconditional defection provides only
one test of the incentives generated by a reputation system.

Behavioral timing changes this test even when average exploitation
is held fixed. In the representative configurations with iid invasion
windows, sufficiently persistent exploitative episodes closed those
windows. Dispersed negative assessments are often interrupted by
positive assessments that reset a warning, whereas persistent
episodes make successive negative assessments more likely.
The resulting difference in reputation changes the benefits received
by the mutant without changing its average saving in cooperation
costs. This effect requires retaining the association between
behavioral mode and internal reputation state. An assessment
probability averaged over modes generally discards that association.
The benchmark comparisons reinforce the importance of the update
protocol: quantitative scores also respond to temporal persistence,
whereas aggregation from the specified stationary infinite pool
removes the temporal ordering of observations. These comparisons
identify consequences of particular processing rules rather than
a general ranking of reputation mechanisms.

Intermediate-state readout provides an additional way to change
these incentives. When a warning begins to reduce the probability
of being treated as good, unfavorable evidence has consequences
before the internal record reaches bad standing. In the representative
cases, this adjustment closed the iid invasion windows while
retaining high resident cooperation. Its scope is nevertheless
limited by the feedback that also sustains the resident state.
Lowering the warning-state readout reduces cooperation and, under
Shunning and Stern Judging, eventually removes the high-reputation
stationary branch. The evolutionary boundary for excluding mutants
and the existence boundary of that branch are mathematically distinct.
An effective readout must therefore provide sufficient sensitivity
to behavioral deviations while preserving the cooperative
environment in which those incentives operate.

Several assumptions delimit these conclusions. The analytical
description uses a mean-field closure for private assessments and
assumes that reputations equilibrate faster than strategy frequencies.
The finite-population simulations examine reputation formation at
fixed behavioral composition; they do not establish the full
evolutionary dynamics of finite populations. Where multiple
reputation branches coexist, our evolutionary comparisons use the
branch selected from initially good records. Basin averages also
refer to a specified distribution of initial strategy compositions.
The representative depths depend on operational recognition
criteria, and equal total depth compares numbers of internal states
without measuring their cognitive costs. Finally, the assessment
norms and updating parameters remain fixed within each evolutionary
comparison. Resistance to the tested action mutants does not
establish stability against arbitrary strategies, competing
assessment rules or alternative memory mechanisms.

The model offers a way to formulate behavioral questions about
reputation revision. Experiments on moral learning and trust repair
show that prior judgments can affect responses to subsequent
behavior, although they do not establish the particular transition
rules studied here \cite{siegel2018,schweitzer2006}. Our results
suggest distinguishing an individual's current social standing
from evidence that has accumulated without yet reversing it.
They also suggest testing whether the same frequency of favorable
and unfavorable observations produces different judgments when
their order changes, and whether those differences depend on
initial standing. Such observations would help identify which
features of directional persistence describe human judgments.

Taken together, our findings connect the retention of social
evidence to the incentives supporting indirect reciprocity.
Preserving reputations can protect useful behavioral distinctions,
but its cooperative value depends on which judgments persist,
how conflicting evidence accumulates and when that evidence affects
subsequent interactions. Understanding how reputations are revised
is therefore part of understanding how social norms sustain
cooperation.

\section*{Methods}\label{sec:methods}

\subsection*{Interactions and private assessment}

We consider an effectively infinite, well-mixed population. In each interaction, one individual acts as donor, another as recipient, and an observer assesses the donor. The donor first consults its private reputation of the recipient and chooses an intended action. Individuals using $\mathrm{ALLC}$ intend to cooperate with every recipient, individuals using $\mathrm{ALLD}$ always intend to defect, and individuals using $\mathrm{DISC}$ intend to cooperate with recipients they regard as good and to defect against recipients they regard as bad. Successful cooperation gives the recipient a benefit $b$ and costs the donor $c$, where $b>c>0$.

The observer assesses the donor's realized action using its own private reputation of the recipient. We represent the assessment rule in second-order form. Cooperation with a good recipient receives a positive assessment and defection against a good recipient receives a negative assessment. The assessment of actions towards a bad recipient is specified by $p$ and $q$, where $p$ is the probability of a positive assessment after cooperation and $q$ is the corresponding probability after defection. The four norms considered in the analyses are Scoring, Shunning, Stern Judging and Simple Standing, corresponding to $(p,q)=(1,0)$, $(0,0)$, $(0,1)$ and $(1,1)$, respectively.

Errors can occur during action implementation and assessment. Intended cooperation fails and is implemented as defection with probability $e_1$, whereas intended defection is implemented without error. After the social norm has been applied to the realized action, the resulting assessment is reversed with probability $e_2$. We assume $0\leq e_1,e_2<1/2$, and the two errors occur independently. Let $P^+_{A\mid R}$ denote the probability that the donor receives a positive assessment, conditional on its intended action $A\in\{\mathrm{C},\mathrm{D}\}$ and the observer's private reputation $R\in\{\mathrm{G},\mathrm{B}\}$ of the recipient. Defining $\epsilon=(1-e_1)(1-e_2)+e_1e_2$, the four effective assessment probabilities are
\begin{equation}
\begin{aligned}
P^+_{\mathrm{C}\mid\mathrm{G}}&=\epsilon,\\
P^+_{\mathrm{D}\mid\mathrm{G}}&=e_2,\\
P^+_{\mathrm{C}\mid\mathrm{B}}
&=p(\epsilon-e_2)+q(1-\epsilon-e_2)+e_2,\\
P^+_{\mathrm{D}\mid\mathrm{B}}
&=q(1-2e_2)+e_2.
\end{aligned}
\label{eq:methods-effective-assessment}
\end{equation}
These expressions integrate over implementation and assessment errors while retaining the distinction between the donor's intended action and the observer's private view of the recipient. Their event-by-event derivation and values under each norm are given in Supplementary Note~1 and Supplementary Table~\ref{tab:si-effective-assessment}.

\subsection*{Inertial reputation updating}

For each observer--target pair $(i,j)$, the observer retains an internal reputation state $L_{ij}$. In the one-step model analysed in the main text, $L_{ij}\in\{\mathrm{G},\mathrm{W},\mathrm{B},\mathrm{S}\}$. A negative assessment moves $\mathrm{G}$ to the warning state $\mathrm{W}$. A further negative assessment moves $\mathrm{W}$ to $\mathrm{B}$, whereas a positive assessment restores $\mathrm{G}$. Conversely, a positive assessment moves $\mathrm{B}$ to the suspicious state $\mathrm{S}$. A further positive assessment moves $\mathrm{S}$ to $\mathrm{G}$, whereas a negative assessment restores $\mathrm{B}$. Positive assessments leave $\mathrm{G}$ unchanged, and negative assessments leave $\mathrm{B}$ unchanged.

Behavior is based on a binary readout of this internal state. States $\mathrm{G}$ and $\mathrm{B}$ are read as good and bad, respectively. State $\mathrm{W}$ is read as good with probability $\alpha$, and state $\mathrm{S}$ is read as good with probability $1-\beta$, where $0\leq\alpha,\beta\leq1$. The internal state persists between assessment events, whereas its binary readout is sampled independently whenever the record is consulted. Thus, intermediate states retain information about the direction from which the current reputation was reached without requiring behavioral strategies or social norms to operate on more than two reputational categories.

For a fixed assessment environment, let $\varphi$ be the probability that a focal individual receives a positive assessment. The internal state then follows a finite Markov chain whose transition matrix is the probability-weighted mixture of the transitions following positive and negative assessments. Let $\pi_{\ell}(\varphi)$ denote its stationary probability of state $\ell$. The probability that the individual is read as good is
\begin{equation}
\mathcal R(\varphi)
=\pi_{\mathrm{G}}(\varphi)
+\alpha\pi_{\mathrm{W}}(\varphi)
+(1-\beta)\pi_{\mathrm{S}}(\varphi).
\label{eq:methods-reputation-map-definition}
\end{equation}
For one-step inertia, Eq.~\eqref{eq:methods-reputation-map-definition} becomes
\begin{equation}
\mathcal R(\varphi)
=\frac{
\varphi^2\left[1+\alpha(1-\varphi)\right]
+\varphi(1-\varphi)^2(1-\beta)
}{1-\varphi+\varphi^2}.
\label{eq:methods-reputation-map}
\end{equation}
When $\alpha=\beta=0$, every positive assessment produces a good binary readout and every negative assessment a bad readout, so that $\mathcal R(\varphi)=\varphi$ and the model recovers instantaneous binary updating. Supplementary Note~1 extends the mechanism to $k_1$ warning states and $k_2$ suspicious states. The corresponding stationary state probabilities and general reputation mapping are derived in Supplementary Eqs.~\eqref{eq:si-stationary-distribution} and \eqref{eq:si-reputation-transform}.

\subsection*{Stationary reputation self-consistency}

Let $\mathcal S=\{\mathrm{ALLC},\mathrm{ALLD},\mathrm{DISC}\}$, and let $f_I$ and $r_I$ denote the population frequency and good-reputation probability of strategy $I\in\mathcal S$, respectively. The average good-reputation probability is $\bar r=\sum_{I\in\mathcal S}f_Ir_I$. For interactions involving a $\mathrm{DISC}$ donor, we use the standard mean-field closure that the donor's and observer's private reputations of the recipient are conditionally independent given the recipient's strategy. The probabilities of agreement on good, agreement on bad and either ordered disagreement are therefore
\[
h_{\mathrm{GG}}=\sum_{I\in\mathcal S}f_Ir_I^2,
\qquad
h_{\mathrm{BB}}=\sum_{I\in\mathcal S}f_I(1-r_I)^2,
\qquad
h_{\mathrm{GB}}=\sum_{I\in\mathcal S}f_Ir_I(1-r_I).
\]

Let $\varphi_I$ be the probability that a type-$I$ donor receives a positive assessment. Combining the strategy's intended action with Eq.~\eqref{eq:methods-effective-assessment} gives
\begin{equation}
\begin{aligned}
\varphi_{\mathrm{ALLC}}
&=\bar rP^+_{\mathrm{C}\mid\mathrm{G}}
+(1-\bar r)P^+_{\mathrm{C}\mid\mathrm{B}},\\
\varphi_{\mathrm{ALLD}}
&=\bar rP^+_{\mathrm{D}\mid\mathrm{G}}
+(1-\bar r)P^+_{\mathrm{D}\mid\mathrm{B}},\\
\varphi_{\mathrm{DISC}}
&=h_{\mathrm{GG}}P^+_{\mathrm{C}\mid\mathrm{G}}
+h_{\mathrm{GB}}
\left(P^+_{\mathrm{C}\mid\mathrm{B}}
+P^+_{\mathrm{D}\mid\mathrm{G}}\right)
+h_{\mathrm{BB}}P^+_{\mathrm{D}\mid\mathrm{B}}.
\end{aligned}
\label{eq:methods-positive-assessment}
\end{equation}
The three terms for $\mathrm{DISC}$ distinguish agreement that the recipient is good, the two possible directions of disagreement, and agreement that the recipient is bad. A detailed event decomposition is provided in Supplementary Eq.~\eqref{eq:si-positive-assessment-probabilities}.

Equations~\eqref{eq:methods-reputation-map} and \eqref{eq:methods-positive-assessment} define a self-consistency map between current reputations and the reputations generated by the resulting assessment environment. Stationary reputations satisfy $r_I^*=\mathcal R(\varphi_I^*)$, where $\mathcal R$ is the stationary readout of the underlying multistate process. To locate a self-consistent solution from the all-good initial condition, we use the auxiliary relaxation
\begin{equation}
\frac{\mathrm dr_I}{\mathrm d\tau}
=\mathcal R(\varphi_I)-r_I,
\qquad I\in\mathcal S.
\label{eq:methods-reputation-dynamics}
\end{equation}
Here $\tau$ is an auxiliary time. Equation~\eqref{eq:methods-reputation-dynamics} has the stationary self-consistency solutions as its fixed points; it is not the microscopic transient dynamics of the internal reputation states. We assume that all individuals are initially regarded as good and set $r_I(0)=1$, corresponding to all internal records in state $\mathrm G$. We compare the selected solution with the asynchronous full-state dynamics. Their stationary equations are equivalent within the same private-assessment closure. Agreement of the selected branch is checked numerically rather than assumed from this equivalence. Full-state equations, the consistency proof and validation procedures are given in Supplementary Note~\ref{sec:si-full-state-dynamics}. Stability and basins inferred from the reduced flow refer to that auxiliary flow unless independently verified with the full-state system.

\subsection*{Payoffs and strategy evolution}

We assume that reputation dynamics are fast relative to changes in strategy frequencies. Strategy composition is therefore held fixed while Eq.~\eqref{eq:methods-reputation-dynamics} relaxes from the all-good initial condition to a stable reputation equilibrium. The equilibrium reputations $r_I^*$ and their population average $\bar r^*=\sum_{I\in\mathcal S}f_Ir_I^*$ then determine payoffs on the slower evolutionary timescale. A type-$I$ recipient receives intended cooperation with probability
\begin{equation}
\kappa_I^{\mathrm{in}}
=f_{\mathrm{ALLC}}+f_{\mathrm{DISC}}r_I^*.
\label{eq:methods-incoming-cooperation}
\end{equation}
The corresponding probabilities of intending cooperation as a donor are one for $\mathrm{ALLC}$, zero for $\mathrm{ALLD}$ and $\bar r^*$ for $\mathrm{DISC}$. Denoting these probabilities by $\kappa_I^{\mathrm{out}}$, the expected payoff is
\begin{equation}
\Pi_I=(1-e_1)
\left(b\kappa_I^{\mathrm{in}}-c\kappa_I^{\mathrm{out}}\right).
\label{eq:methods-payoff}
\end{equation}
The factor $1-e_1$ is common to all strategies and rescales evolutionary time, but implementation error also changes $r_I^*$ through the assessment environment. The explicit payoffs are given in Supplementary Note~\ref{sec:si-model}.

The equilibrium payoffs determine strategy frequencies through the replicator equations
\begin{equation}
\frac{\mathrm df_I}{\mathrm dt}
=f_I(\Pi_I-\bar\Pi).
\label{eq:methods-replicator-dynamics}
\end{equation}
Here $\bar\Pi=\sum_{J\in\mathcal S}f_J\Pi_J$. At every strategy composition visited by Eq.~\eqref{eq:methods-replicator-dynamics}, the reputation equilibrium is recomputed before the payoffs are evaluated. We distinguish the rate at which cooperation is intended,
\begin{equation}
\rho_{\mathrm{int}}
=f_{\mathrm{ALLC}}+f_{\mathrm{DISC}}\bar r^*,
\label{eq:methods-intended-cooperation-rate}
\end{equation}
from the rate at which cooperation is realized,
\begin{equation}
\rho_{\mathrm{C}}
=(1-e_1)\rho_{\mathrm{int}}.
\label{eq:methods-cooperation-rate}
\end{equation}
Intended cooperation, $\rho_{\mathrm{int}}$, and realized cooperation, $\rho_{\mathrm C}=(1-e_1)\rho_{\mathrm{int}}$, are reported separately. The equilibrium percentages in Fig.~\ref{fig:strategy-dynamics} and Extended Data Fig.~\ref{fig:ed-strategy-alpha-beta}, and the resident cooperation curves in Extended Data Fig.~\ref{fig:ed-alpha-invasion}, describe intended cooperation. The depth and benefit-to-cost scans in Fig.~\ref{fig:strategy-depth} and Extended Data Figs.~\ref{fig:ed-strategy-bc} and~\ref{fig:ed-strategy-depth-all} retain the realized-cooperation scale of the supplied figures. We use $M=512$ common initial strategy compositions sampled uniformly from the simplex. Their basin-weighted averages satisfy
\begin{equation}
\left\langle\rho_{\mathrm C}\right\rangle_{\Delta}
=(1-e_1)\left\langle\rho_{\mathrm{int}}\right\rangle_{\Delta}
=\frac{1-e_1}{M}\sum_{m=1}^{M}\rho_{\mathrm{int}}\!\left(\boldsymbol f_{\mathrm{end}}^{(m)}\right).
\label{eq:methods-basin-weighted-cooperation}
\end{equation}
Here $\boldsymbol f_{\mathrm{end}}^{(m)}$ is the terminal strategy composition. A terminal estimate is interpreted as an equilibrium outcome only after its convergence checks pass. At fixed $e_1$, the common factor $1-e_1$ preserves comparisons between allocations. This rescaling does not automatically convert a basin defined by a cooperation cutoff. In Fig.~\ref{fig:strategy-dynamics}, shading is evaluated directly from intended cooperation of at least $0.5$. The three unconfirmed $(1,1)$ depth-scan estimates identified in the figure captions are retained as provisional terminal estimates.

\subsection*{Numerical integration}

We integrated the auxiliary reputation and replicator equations in Python using SciPy's adaptive \texttt{solve\_ivp} implementation with the explicit DOP853 method. For every fixed strategy composition, reputations were initialized at $r_I(0)=1$ and integrated with relative tolerance $10^{-9}$, absolute tolerance $10^{-11}$ and maximum auxiliary time $2\times10^4$. A reputation solution was accepted only when $\max_I|\mathrm dr_I/\mathrm d\tau|\leq10^{-10}$. Strategy trajectories were integrated with relative tolerance $10^{-8}$, absolute tolerance $10^{-10}$, maximum step 20 and maximum time $2\times10^4$; convergence required $\max_I|\mathrm df_I/\mathrm dt|\leq10^{-7}$. The equilibrated strategy vector field was evaluated on a triangular simplex lattice and interpolated piecewise linearly for trajectory integration, after which reputations and cooperation were recomputed at every terminal composition without interpolation. Runs that reached the maximum time without satisfying the residual criterion were flagged as non-converged rather than treated as equilibria. Small regression tests against the fixed-step implementation and against the implicit Radau method are reported in Supplementary Note~\ref{sec:si-evolution}.

\subsection*{Local evolutionary stability of discriminators}

We derived the local stability boundaries of a pure $\mathrm{DISC}$ population by introducing either $\mathrm{ALLD}$ or $\mathrm{ALLC}$ at vanishing frequency. Let $r_{\mathrm R}^{*}$ denote the resident reputation selected from the all-good initial condition, and let $r_{\mathrm D}^{*}$ and $r_{\mathrm C}^{*}$ be the equilibrium reputations of rare $\mathrm{ALLD}$ and $\mathrm{ALLC}$ mutants in this resident environment. Because a rare mutant does not alter the resident environment, its reputation is obtained by applying $\mathcal R$ to its assessment probability against resident recipients. Equating the payoff of a rare $\mathrm{ALLD}$ mutant with that of the resident gives
\begin{equation}
R_{\mathrm L}=\frac{r_{\mathrm R}^{*}}{r_{\mathrm R}^{*}-r_{\mathrm D}^{*}}.
\label{eq:methods-lower-benefit-cost-threshold}
\end{equation}
Equating the payoff of a rare $\mathrm{ALLC}$ mutant with that of the resident gives
\begin{equation}
R_{\mathrm U}=\frac{1-r_{\mathrm R}^{*}}{r_{\mathrm C}^{*}-r_{\mathrm R}^{*}}.
\label{eq:methods-upper-benefit-cost-threshold}
\end{equation}
When both denominators are positive, pure $\mathrm{DISC}$ resists both mutants if $R_{\mathrm L}<b/c<R_{\mathrm U}$. If $r_{\mathrm R}^{*}\leq r_{\mathrm D}^{*}$, no finite $b/c$ excludes $\mathrm{ALLD}$. If $r_{\mathrm C}^{*}\leq r_{\mathrm R}^{*}$, $\mathrm{ALLC}$ imposes no finite upper threshold. We used the signs of the two mutant growth rates to classify the $e_1$--$e_2$, $\alpha$--$\beta$ and $(k_1,k_2)$ parameter planes. These local invasion tests are distinct from stability within an invariant simplex edge and from the sizes of attraction basins. The expanded reputational expressions, invasion exponents and limiting cases are derived in Supplementary Note~\ref{sec:si-evolutionary-stability}.

\subsection*{Rare partial and temporally persistent mutants}
\label{sec:methods-opportunism}

For the invasion analyses we hold the population at the all-good-selected pure-DISC reputation equilibrium. A rare under-cooperation mutant uses DISC, except that each intended cooperation is independently replaced by intended defection with probability $\delta$. This choice precedes the ordinary implementation error. In the resident environment, let $\Phi_{\mathrm C}^*$ and $\Phi_{\mathrm D}^*$ be the mutant's positive-assessment probabilities conditional on intending cooperation and defection, and let $\Delta\Phi^*=\Phi_{\mathrm C}^*-\Phi_{\mathrm D}^*$. Its iid assessment probability is $\Phi_{\mathrm X}(\delta)=\Phi_{\mathrm R}^*-\delta r_{\mathrm R}^*\Delta\Phi^*$, and its stationary good-reputation probability is $r_{\mathrm X}^*(\delta)=\mathcal R[\Phi_{\mathrm X}(\delta)]$. The normalized invasion exponent is Eq.~\eqref{eq:results-partial-invasion}. We classify a parameter setting as secure against this tested family only when no $\delta\in(0,1]$ has a positive exponent. The local derivative at $\delta=0$ and the full interval scan answer different questions.

For temporal mutants, a behavioural mode $Z\in\{0,1\}$ uses DISC when $Z=0$ and ALLD when $Z=1$. In donor-action time, mode transitions after the action according to $\Pr(0\to1)=\delta(1-\eta_{\mathrm a})$ and $\Pr(1\to0)=(1-\delta)(1-\eta_{\mathrm a})$. The stationary exploitative fraction is $\delta$. At every donor action the mutant forms an intention, implementation error determines the realized action and its payoff, and a focal observer samples that action independently with probability $p_{\mathrm{obs}}$. If sampled, the observer applies the social norm and assessment error and updates its latent record; the mutant then switches mode, whether or not the focal observer sampled the action. Readouts are sampled whenever a record is consulted. For one rare mutant among $N=150$ individuals and an observer sampled from the other $N-1$ DISC residents, $p_{\mathrm{obs}}=1/149$. Consecutive observations see persistence
\begin{equation}
\eta=\frac{p_{\mathrm{obs}}\eta_{\mathrm a}}{1-(1-p_{\mathrm{obs}})\eta_{\mathrm a}},
\qquad 0\leq\eta<1.
\label{eq:methods-observed-persistence}
\end{equation}
Figure~\ref{fig:strategy-invasion} and Extended Data Fig.~\ref{fig:ed-aggregation-quantitative} use $\eta$ as their temporal coordinate. We obtain the stationary mutant reputation from the joint chain of behavioural mode and latent reputation state, then evaluate Eq.~\eqref{eq:results-partial-invasion} with $r_{\mathrm X}^*(\delta,\eta)$. Observation thinning does not alter the stationary resident reputation; it changes the correlation seen by a focal record. The joint-chain construction, endpoint limits and benchmark-specific observation protocols are given in Supplementary Note~\ref{sec:si-evolutionary-stability}.

\subsection*{Readout security and resident branch limits}
\label{sec:methods-readout-security}

For each $(\alpha,\beta)$ we recomputed the selected pure-DISC resident reputation, holding the norm, memory structure and error parameters fixed. Write $r^*_{\alpha,\beta}$ for this reputation, $u^*_{\alpha,\beta}=\Phi_{\mathrm R}(r^*_{\alpha,\beta})$ for its positive-assessment probability, and $\Delta\Phi^*_{\alpha,\beta}=\Phi_{\mathrm C}(r^*_{\alpha,\beta})-\Phi_{\mathrm D}(r^*_{\alpha,\beta})$. Define the reputation penalty per unit saved cooperation cost by
\begin{equation}
\begin{aligned}
g_{\alpha,\beta}(\delta)
&=\frac{r^*_{\alpha,\beta}-\mathcal R_{\alpha,\beta}
\left(u^*_{\alpha,\beta}-\delta r^*_{\alpha,\beta}\Delta\Phi^*_{\alpha,\beta}\right)}
{\delta r^*_{\alpha,\beta}},\qquad \delta>0,\\
g_{\alpha,\beta}(0)
&=\Delta\Phi^*_{\alpha,\beta}\mathcal R'_{\alpha,\beta}(u^*_{\alpha,\beta}).
\end{aligned}
\label{eq:methods-readout-penalty}
\end{equation}
Here the prime denotes differentiation with respect to the assessment probability. Complete exclusion of the independent under-cooperation family is equivalent to $(b/c)\min_{0\leq\delta\leq1}g_{\alpha,\beta}(\delta)\geq1$. The warning-readout boundary $\alpha_{\mathrm{sec}}$ is the largest readout satisfying this condition on the specified resident branch. We evaluated the analytical limit at $\delta=0$, the endpoint $\delta=1$ and a grid of 2,001 exploitation probabilities, refining detected interior minima with bounded scalar optimization and refining readout boundaries by root finding. For $k_2=0$ and $\Delta\Phi^*_{\alpha,\beta}>0$, concavity of $\mathcal R$ makes the local and complete-family conditions equivalent. For configurations with recovery states, we checked the complete criterion separately. An endpoint value $\alpha_{\mathrm{sec}}=1$ denotes security up to the parameter-domain boundary, not necessarily a root of the local equality.

The separate high-reputation branch limit $\alpha_{\mathrm{rep}}$ was located as a non-degenerate fold of $F(r,\alpha,\beta)=\mathcal R_{\alpha,\beta}[\Phi_{\mathrm R}(r)]-r$, solving $F=\partial_rF=0$ at fixed $\beta$. This stationary condition contains no payoff parameter. It is distinct from the invasion criterion and from stability of the complete latent-state dynamics, assessed as described in Supplementary Note~\ref{sec:si-full-state-dynamics}. We retained the different high- and lower-branch outcomes when interpreting Extended Data Fig.~\ref{fig:ed-alpha-invasion}. The mapping derivatives, exact error-free security threshold, finite-error boundaries and recovery-readout comparison are given in Supplementary Note~\ref{sec:si-readout-security}.

\subsection*{Finite-population Monte Carlo simulations}
\label{sec:methods-finite-monte-carlo}

We complemented the population-level reputation dynamics with an individual-based model that retains the private internal state of each reputation record. Unless otherwise stated, the simulated population contained $N=150$ individuals, with equal numbers of $\mathrm{ALLC}$, $\mathrm{ALLD}$ and $\mathrm{DISC}$. Strategy frequencies were held fixed: these simulations measure reputation formation within a given population rather than strategy evolution. We stored the records held by $\mathrm{DISC}$ observers because only these records can affect behaviour. Their targets nevertheless comprised all three strategies, allowing us to ask whether discriminators distinguish conditional cooperators and unconditional cooperators from unconditional defectors.

Each Monte Carlo event drew a donor and a distinct recipient uniformly from the population, followed independently by an observer drawn uniformly from all $\mathrm{DISC}$ individuals. The observer could coincide with either interacting individual. Self-images were retained and included in the estimators. The donor selected an intended action from its strategy and, if it used $\mathrm{DISC}$, from a newly sampled binary readout of its internal record of the recipient. After implementation error, the observer evaluated the realized action using an independently sampled readout of its own record of the recipient. The social norm and assessment error then determined whether the observer's internal record of the donor received a positive or negative
update. Thus, every consultation resampled the binary readout according to Supplementary Eq.~\eqref{eq:si-binary-readout}, whereas the internal state $L_{ij}$ persisted between events and followed the same inertial transitions as in Supplementary Fig.~\ref{fig:si-state-transitions}.

One Monte Carlo sweep comprised $n_{\mathrm{DISC}}N$ events, equal to the number of retained observer--target records and hence to one update event per retained record after averaging across the matrix. Simulations began with all retained records in state $\mathrm{G}$, ran for $1{,}000$ sweeps, and were sampled every five sweeps after a burn-in of $800$ sweeps. For each norm and depth-grid combination we used 30 independent replicates. Endpoint matrix panels are illustrative single-run snapshots and are not used as replicate-level estimates of the quantitative depth statistics. The displayed reputation matrices contain only $\mathrm{DISC}$ observer rows and show separately saved fixed-seed illustrative configurations; all quantitative comparisons use post-burn-in time
averages from every replicate. Uncertainty intervals are percentile 95\% bootstrap intervals across independent replicate means.

For each target strategy $I$, we measured the mean probability $g_I$ that a record is read as good and the probability $d_I$ that two distinct $\mathrm{DISC}$ observers return different binary reputations for the same target. The latter includes stochastic readout of intermediate states and therefore measures disagreement in the behaviourally expressed reputations, not merely variation in latent states. For the equal-frequency mixed populations, reputational separation is
\begin{equation}
\Delta g=\frac{g_{\mathrm{DISC}}+g_{\mathrm{ALLC}}}{2}-g_{\mathrm{ALLD}}.
\label{eq:methods-recognition-separation}
\end{equation}
We average observer disagreement across targets as
\begin{equation}
\bar d=\sum_{I\in\mathcal S}\frac{n_I}{N}d_I,
\label{eq:methods-recognition-disagreement}
\end{equation}
where $n_I$ is the number of type-$I$ targets and $d_I$ averages, over targets of that type, the probability that two distinct DISC observers express different binary images of the same target. Self-images are included as in the simulations. The finite-observer estimators integrate over readout randomness and are given with the event sequence in Supplementary Note~\ref{sec:si-finite-monte-carlo}.

To examine inertia beyond the one-step mechanism, we repeated the same fixed-composition simulations while varying the number of warning states $k_1$, the number of suspicious states $k_2$, or both. All analyses use the direct-reset transitions defined in Supplementary Fig.~\ref{fig:si-state-transitions}: any positive assessment sends a warning directly to $\mathrm G$, and any negative assessment sends a suspicious state directly to $\mathrm B$. Increasing $k_1$ or $k_2$ therefore raises the number of consecutive assessments required to move between established good and established bad reputations without changing the binary action and assessment rules.






\section*{Data availability}
All data generated or analysed during this study are included within the paper and its supplementary information file. 

\section*{Acknowledgements}
xxxx

\section*{Author contributions}
xxxx

\section*{Competing interests}
The authors declare no competing interests.

\begin{thebibliography}{99}
\bibitem{nowak1998}
Nowak, M. A. \& Sigmund, K.
Evolution of indirect reciprocity by image scoring.
\textit{Nature} \textbf{393}, 573--577 (1998).
\url{https://doi.org/10.1038/31225}.

\bibitem{wedekind2000}
Wedekind, C. \& Milinski, M.
Cooperation through image scoring in humans.
\textit{Science} \textbf{288}, 850--852 (2000).
\url{https://doi.org/10.1126/science.288.5467.850}.

\bibitem{schmid2021}
Schmid, L., Chatterjee, K., Hilbe, C. \& Nowak, M. A.
A unified framework of direct and indirect reciprocity.
\textit{Nature Human Behaviour} \textbf{5}, 1292--1302 (2021).
\url{https://doi.org/10.1038/s41562-021-01114-8}.

\bibitem{michelmata2024}
Michel-Mata, S., Kawakatsu, M., Sartini, J., Kessinger, T. A.,
Plotkin, J. B. \& Tarnita, C. E.
The evolution of private reputations in information-abundant landscapes.
\textit{Nature} \textbf{634}, 883--889 (2024).
\url{https://doi.org/10.1038/s41586-024-07977-x}.

\bibitem{effron2010}
Effron, D. A. \& Monin, B.
Letting people off the hook: when do good deeds excuse transgressions?
\textit{Personality and Social Psychology Bulletin} \textbf{36}, 1618--1634 (2010).
\url{https://doi.org/10.1177/0146167210385922}.

\bibitem{siegel2018}
Siegel, J. Z., Mathys, C., Rutledge, R. B. \& Crockett, M. J.
Beliefs about bad people are volatile.
\textit{Nature Human Behaviour} \textbf{2}, 750--756 (2018).
\url{https://doi.org/10.1038/s41562-018-0425-1}.

\bibitem{schweitzer2006}
Schweitzer, M. E., Hershey, J. C. \& Bradlow, E. T.
Promises and lies: restoring violated trust.
\textit{Organizational Behavior and Human Decision Processes} \textbf{101}, 1--19 (2006).
\url{https://doi.org/10.1016/j.obhdp.2006.05.005}.

\bibitem{nowak2005}
Nowak, M. A. \& Sigmund, K.
Evolution of indirect reciprocity.
\textit{Nature} \textbf{437}, 1291--1298 (2005).
\url{https://doi.org/10.1038/nature04131}.

\bibitem{ohtsuki2004}
Ohtsuki, H. \& Iwasa, Y.
How should we define goodness?---Reputation dynamics in indirect reciprocity.
\textit{Journal of Theoretical Biology} \textbf{231}, 107--120 (2004).
\url{https://doi.org/10.1016/j.jtbi.2004.06.005}.

\bibitem{hilbe2018}
Hilbe, C., Schmid, L., Tkadlec, J., Chatterjee, K. \& Nowak, M. A.
Indirect reciprocity with private, noisy, and incomplete information.
\textit{Proceedings of the National Academy of Sciences USA} \textbf{115}, 12241--12246 (2018).
\url{https://doi.org/10.1073/pnas.1810565115}.

\bibitem{pacheco2006}
Pacheco, J. M., Santos, F. C. \& Chalub, F. A. C. C.
Stern-judging: a simple, successful norm which promotes cooperation
under indirect reciprocity.
\textit{PLoS Computational Biology} \textbf{2}, e178 (2006).
\url{https://doi.org/10.1371/journal.pcbi.0020178}.

\bibitem{ohtsuki2007}
Ohtsuki, H. \& Iwasa, Y.
Global analyses of evolutionary dynamics and exhaustive search for
social norms that maintain cooperation by reputation.
\textit{Journal of Theoretical Biology} \textbf{244}, 518--531 (2007).
\url{https://doi.org/10.1016/j.jtbi.2006.08.018}.

\bibitem{brandt2004}
Brandt, H. \& Sigmund, K.
The logic of reprobation: assessment and action rules for indirect
reciprocation.
\textit{Journal of Theoretical Biology} \textbf{231}, 475--486 (2004).
\url{https://doi.org/10.1016/j.jtbi.2004.06.032}.

\bibitem{ohtsuki2006}
Ohtsuki, H. \& Iwasa, Y.
The leading eight: social norms that can maintain cooperation by indirect reciprocity.
\textit{Journal of Theoretical Biology} \textbf{239}, 435--444 (2006).
\url{https://doi.org/10.1016/j.jtbi.2005.08.008}.

\bibitem{fujimoto2022}
Fujimoto, Y. \& Ohtsuki, H.
Reputation structure in indirect reciprocity under noisy and private assessment.
\textit{Scientific Reports} \textbf{12}, 10500 (2022).
\url{https://doi.org/10.1038/s41598-022-14171-4}.

\bibitem{kessinger2023}
Kessinger, T. A., Tarnita, C. E. \& Plotkin, J. B.
Evolution of norms for judging social behavior.
\textit{Proceedings of the National Academy of Sciences USA}
\textbf{120}, e2219480120 (2023).
\url{https://doi.org/10.1073/pnas.2219480120}.

\bibitem{radzvilavicius2021}
Radzvilavicius, A. L., Kessinger, T. A. \& Plotkin, J. B.
Adherence to public institutions that foster cooperation.
\textit{Nature Communications} \textbf{12}, 3567 (2021).
\url{https://doi.org/10.1038/s41467-021-23783-9}.

\bibitem{kawakatsu2024}
Kawakatsu, M., Kessinger, T. A. \& Plotkin, J. B.
A mechanistic model of gossip, reputations, and cooperation.
\textit{Proceedings of the National Academy of Sciences USA} \textbf{121}, e2400689121 (2024).
\url{https://doi.org/10.1073/pnas.2400689121}.

\bibitem{radzvilavicius2019}
Radzvilavicius, A. L., Stewart, A. J. \& Plotkin, J. B.
Evolution of empathetic moral evaluation.
\textit{eLife} \textbf{8}, e44269 (2019).
\url{https://doi.org/10.7554/eLife.44269}.

\bibitem{tham2026}
Tham, Y. J., Hilbe, C. \& Murase, Y.
Indirect reciprocity with dual private assessment.
\textit{Proceedings of the National Academy of Sciences USA} \textbf{123}, e2624656123 (2026).
\url{https://doi.org/10.1073/pnas.2624656123}.


\bibitem{tanabe2013}
Tanabe, S., Suzuki, H. \& Masuda, N.
Indirect reciprocity with trinary reputations.
\textit{Journal of Theoretical Biology} \textbf{317}, 338--347 (2013).
\url{https://doi.org/10.1016/j.jtbi.2012.10.031}.

\bibitem{murase2022}
Murase, Y., Kim, M. \& Baek, S. K.
Social norms in indirect reciprocity with ternary reputations.
\textit{Scientific Reports} \textbf{12}, 455 (2022).
\url{https://doi.org/10.1038/s41598-021-04033-w}.

\bibitem{schmid2023}
Schmid, L., Ekbatani, F., Hilbe, C. \& Chatterjee, K.
Quantitative assessment can stabilize indirect reciprocity under imperfect information.
\textit{Nature Communications} \textbf{14}, 2086 (2023).
\url{https://doi.org/10.1038/s41467-023-37817-x}.



% \bibitem{sasaki2017}
% Sasaki, T., Okada, I. \& Nakai, Y.
% The evolution of conditional moral assessment in indirect reciprocity.
% \textit{Scientific Reports} \textbf{7}, 41870 (2017).
% \url{https://doi.org/10.1038/srep41870}.

\end{thebibliography}



% Extended Data fragment; labels match the Results and existing Methods/SI.
\clearpage
\begingroup
\setcounter{figure}{0}
\renewcommand{\figurename}{Extended Data Fig.}
\ifdefined\theHfigure
\renewcommand{\theHfigure}{ED\arabic{figure}}
\fi

\begin{figure}[p]
\centering
\includegraphics[width=0.92\linewidth,height=0.70\textheight,keepaspectratio]{ED_Fig1_bistability.pdf}
\caption{\textbf{The range of one-step bistability depends on the social norm.}
\textbf{a--d}, Stationary good-reputation probability in homogeneous DISC populations across tolerance $\alpha$ and strictness $\beta$, under Scoring, Shunning, Stern Judging and Simple Standing. Colour shows the solution selected from $r_{\mathrm{DISC}}(0)=1$ using Methods, Eq.~\eqref{eq:methods-reputation-dynamics}. Hatching marks regions with two attracting roots separated by an unstable root of that equilibration equation. Squares mark instantaneous updating, $(\alpha,\beta)=(0,0)$, and stars mark full one-step inertia, $(1,1)$. The internal paths remain fixed at $k_1=k_2=1$, and $e_1=e_2=0.02$. The relation between stationary self-consistency and full-state reputation dynamics is given in Supplementary Note~\ref{sec:si-full-state-dynamics}.}
\label{fig:ed-bistability-alpha-beta}
\end{figure}

\begin{figure}[p]
\centering
\includegraphics[width=\linewidth,height=0.72\textheight,keepaspectratio]{ED_Fig2_Finite-population reputation dynamics.pdf}
\caption{\textbf{Strategy-specific reputations and observer disagreement across inertia depths.}
\textbf{a--d}, Rows show Scoring, Shunning, Stern Judging and Simple Standing. The first three columns give the mean probabilities that DISC observers read DISC, ALLC and ALLD targets as good; the last column gives target-averaged disagreement $\bar d$ between distinct observers. Each panel varies tolerance depth $k_1$ and strictness depth $k_2$ at full readout, $\alpha=\beta=1$. Black and red outlines mark the binary baseline $(0,0)$ and full one-step inertia $(1,1)$. Values are means of 30 independent simulations with $N=150$, 50 individuals per strategy, initially good records and $e_1=e_2=0.02$. Strategies remain fixed. Estimators and the extended sampling windows used for slowly relaxing cases are specified in Methods and Supplementary Note~\ref{sec:si-selection-robustness}.}
\label{fig:ed-finite-population-reputation-dynamics}
\end{figure}

\begin{figure}[p]
\centering
\includegraphics[width=\linewidth]{ED_Fig3_recognition selection.pdf}
\caption{\textbf{Representative memory configurations depend on the recognition criteria.}
Columns show Scoring, Shunning, Stern Judging and Simple Standing.
\textbf{a--d}, Mean reputational separation against mean observer disagreement for configurations from the $0\leq k_1,k_2\leq6$ scan; the displayed axes focus on the region of relatively high recognition. Colour denotes total depth $K=k_1+k_2$. Open circles mark configurations that are non-dominated when smaller total depth, greater separation and lower disagreement are considered jointly. Stars mark the representative shortest pairs selected using $\Delta g\geq0.8$ and $\bar d\leq0.1$, indicated by dashed lines; error bars show pointwise 95\% bootstrap intervals.
\textbf{e--h}, Selected pairs when the minimum separation and maximum disagreement thresholds are varied. Colour denotes their total depth, and outlined cells mark the default criteria. Selection first minimizes total depth, then resolves ties by greater separation and lower disagreement. Each configuration uses 30 independent runs with $N=150$, equal strategy frequencies, $\alpha=\beta=1$ and $e_1=e_2=0.02$. These thresholds identify representative configurations within the surveyed grid.}
\label{fig:ed-selection-robustness}
\end{figure}

\begin{figure}[p]
\centering
\includegraphics[width=0.96\linewidth,height=0.69\textheight,keepaspectratio]{ED_Fig4_strategy dynamics.pdf}
\caption{\textbf{The two directions of one-step inertia have different evolutionary effects.}
\textbf{a--d}, Strategy dynamics under Scoring, Shunning, Stern Judging and Simple Standing. Upper and lower rows use $(\alpha,\beta)=(1,0)$ and $(0,1)$, respectively, with both internal paths retained at $k_1=k_2=1$. Percentages give intended cooperation at highlighted equilibria. Shading identifies cooperative attraction basins, and filled and open circles denote stable and unstable equilibria; edge classifications refer to stability within the corresponding invariant edge.
\textbf{e--h}, Basin-weighted cooperation across the $(\alpha,\beta)$ plane, averaged over 512 common initial strategy compositions.
\textbf{i--l}, Local invasion outcome of pure DISC. Rose denotes invasion by ALLD, ochre denotes exclusion of ALLD but invasion by ALLC, and green denotes exclusion of both. Dashed and solid contours mark zero growth of ALLD and ALLC. All panels use $b/c=5$, $e_1=e_2=0.02$ and reputation solutions selected from initially good records.}
% The final e--h plotting version must determine whether its stored cooperation
% scale is realized or divided by (1-e_1); two source plotting paths exist.
\label{fig:ed-strategy-alpha-beta}
\end{figure}

\begin{figure}[p]
\centering
\includegraphics[width=\linewidth]{ED_Fig5_bcrule.pdf}
\caption{\textbf{Cooperative outcomes and local invasion boundaries across benefit-to-cost ratios.}
\textbf{a--d}, Basin-weighted realized cooperation under Scoring, Shunning, Stern Judging and Simple Standing. Grey dashed, orange dash-dotted and blue solid curves show the binary baseline, full one-step inertia, and representative full multi-step inertia at $(k_1,k_2)=(2,0),(3,0),(4,3)$ and $(1,2)$, respectively. Each value averages terminal cooperation over 512 common initial strategy compositions. Saturated curve segments satisfy $R_{\mathrm L}<b/c<R_{\mathrm U}$ and indicate resistance of pure DISC to both rare ALLD and rare ALLC. Pale segments lie outside this interval. Downward and upward triangles mark $R_{\mathrm L}$ and $R_{\mathrm U}$ when within the displayed range; labels report upper bounds for the multi-step configurations. The vertical dotted line marks $b/c=5$. These local thresholds concern pure DISC and need not mark the disappearance of cooperative mixed states. Parameters are $e_1=e_2=0.02$; all reputation calculations start from good records.}
\label{fig:ed-strategy-bc}
\end{figure}

\begin{figure}[p]
\centering
\includegraphics[width=\linewidth]{ED_fig6_inertia depth.pdf}
\caption{\textbf{Cooperation and local protection occupy different regions of the depth plane.}
\textbf{a--d}, Basin-weighted realized cooperation across tolerance depth $k_1$ and strictness depth $k_2$ under Scoring, Shunning, Stern Judging and Simple Standing. Colours average terminal cooperation over 512 common initial strategy compositions. Dotted cells identify configurations for which the selected pure-DISC resident excludes both rare ALLD and rare ALLC. Parameters are $\alpha=\beta=1$, $b/c=5$ and $e_1=e_2=0.02$, with initially good reputation records. As in Fig.~\ref{fig:strategy-depth}, the terminal cooperation estimates at $(1,1)$ under Scoring, Shunning and Simple Standing did not pass the stored observable-convergence check and remain provisional.}
\label{fig:ed-strategy-depth-all}
\end{figure}

\begin{figure}[p]
\centering
\includegraphics[width=\linewidth,height=0.68\textheight,keepaspectratio]{ED_Fig7_readout.pdf}
\caption{\textbf{Warning-state readout jointly constrains invasion resistance and resident cooperation.}
Columns show Scoring, Shunning, Stern Judging and Simple Standing, using representative depths $(2,0),(3,0),(4,3)$ and $(1,2)$.
\textbf{a--d}, Width $\delta_{\mathrm c}$ of the positive-invasion set for independently acting under-cooperation mutants against warning-state readout $\alpha$, at $b/c=5$. For the displayed nonempty intervals, the width equals the upper endpoint. Open circles mark the largest secure readout $\alpha_{\mathrm{sec}}$: $0.807064$, $0.806464$, $0.803604$ and $1$, in column order. Solid curves refer to the high branch where it exists; dashed curves in Shunning and Stern Judging show lower-branch invasion below its fold. Security up to $\alpha=1$ in Simple Standing is a parameter-domain endpoint.
\textbf{e--h}, Resident intended cooperation on the corresponding reputation branches. Triangles in Shunning and Stern Judging mark stationary folds at $\alpha_{\mathrm{rep}}=0.591161$ and $0.566986$; dashed curves show the selected lower-branch outcomes below these limits. The fold endpoint itself has a coalesced stationary root. Filled points and vertical dotted lines identify $\alpha=0.8$.
\textbf{i--l}, Security boundaries across $b/c$. Green shading marks exclusion of the tested family on the displayed resident branch; in Shunning and Stern Judging the high-branch secure region lies between $\alpha_{\mathrm{rep}}$ (grey dashed line) and $\alpha_{\mathrm{sec}}$. Grey shading indicates absence of that high branch. At fixed reputation parameters, $\alpha_{\mathrm{rep}}$ is independent of $b/c$, whereas $\alpha_{\mathrm{sec}}$ follows from mutant payoffs. Scoring and Simple Standing have no corresponding fold in this scan, and security alone need not imply high cooperation. Dotted guides mark $b/c=5$ and $\alpha=0.8$. All panels use $\beta=1$ and $e_1=e_2=0.02$. Definitions and derivations are given in Methods and Supplementary Note~\ref{sec:si-readout-security}.}
\label{fig:ed-alpha-invasion}
\end{figure}

\begin{figure}[p]
\centering
\includegraphics[width=\linewidth,height=0.64\textheight,keepaspectratio]{ED_Fig8_mechanism comparison.pdf}
\caption{\textbf{Under-cooperation can exploit tolerance in representative aggregation and quantitative mechanisms.}
All panels use Stern Judging, $b/c=5$, $e_1=e_2=0.02$ and the selected high-reputation pure-DISC resident.
\textbf{a,b}, Normalized invasion exponent for independently acting under-cooperation mutants. The coloured Aggregation curve samples $M=2$ interactions and requires at least $Q=1$ positive assessment ($q=Q/M=1/2$). Its grey parameter family uses $(M,Q)=(2,2),(4,1),(4,2),(4,3),(4,4),(6,3),(8,4),(10,5)$. The coloured Quantitative curve uses an integer score $z\in[-R,R]$ with $R=5$, increments or decrements the score by one after each assessment with truncation at the bounds, and reads $z\geq S=0$ as good. Its grey family uses $R=1,2,3,4$ and $S=0$. Open circles mark positive-invasion interval endpoints for the representative rules.
\textbf{c,d}, Invasion landscapes for those representative rules against mean exploitation $\delta$ and observed persistence $\eta$. Warm and blue colours denote positive and negative invasion exponents. Black curves mark the representative zero-invasion boundaries; grey curves mark zero-invasion boundaries for other parameter choices. Aggregation independently samples with replacement from a stationary infinite interaction pool, yielding persistence-independent outcomes in this specified limit. Quantitative updating retains ordered score changes and is evaluated using a joint chain of behavioural mode and score. The edge $\eta=1$ denotes a one-sided limit. These published processing rules~\cite{michelmata2024,schmid2023} are adapted to the common invasion framework described in Supplementary Note~\ref{sec:si-benchmark-mechanisms}. The plotted quantity is $\Lambda=(\Pi_{\mathrm X}-\Pi_{\mathrm R})/[(1-e_1)c]$.}
\label{fig:ed-aggregation-quantitative}
\end{figure}

\clearpage
\endgroup

\clearpage





\newpage
\pagebreak
\clearpage
\onecolumn
\setcounter{equation}{0}
\setcounter{figure}{0}
\setcounter{table}{0}
\setcounter{section}{0}
% \setcounter{page}{1}

\renewcommand{\thefigure}{\arabic{figure}}
\renewcommand{\figurename}{{\bf Supplementary Figure.}}
\renewcommand{\theequation}{S\arabic{equation}}
\renewcommand{\theHequation}{S\arabic{equation}}

% ==========================================================
% Supplementary Note numbering and heading format
% ==========================================================
\renewcommand{\thesection}{\arabic{section}}
\renewcommand{\thesubsection}{\thesection.\arabic{subsection}}
\renewcommand{\thesubsubsection}{\thesubsection.\arabic{subsubsection}}

\titleformat{\section}
  {\normalfont\Large\bfseries}
  {Supplementary Note~\thesection}
  {0.6em}
  {#1}

\titleformat{\subsection}
  {\normalfont\large\bfseries}
  {Supplementary Note~\thesubsection}
  {0.6em}
  {#1}

\titleformat{\subsubsection}
  {\normalfont\normalsize\bfseries}
  {Supplementary Note~\thesubsubsection}
  {0.6em}
  {#1}

% ==========================================================
% Supplementary Note format in the table of contents
% ==========================================================
\renewcommand{\cftsecpresnum}{Supplementary Note~}
\renewcommand{\cftsubsecpresnum}{Supplementary Note~}
\renewcommand{\cftsubsubsecpresnum}{Supplementary Note~}

\setlength{\cftsecnumwidth}{10em}
\setlength{\cftsubsecnumwidth}{11em}
\setlength{\cftsubsubsecnumwidth}{12em}

\newcounter{fnnumber}
\renewcommand{\thefootnote}{\fnsymbol{footnote}}
\begin{center}
{\Huge ---Supplementary Material---}\\[1em]
\textbf{\huge
Evolutionary dynamics of indirect reciprocity with inertial reputation transitions}\\
\vspace{0.5cm}
Wenqiang Zhu\\
\vspace{0.5cm}
\end{center}


% Show section, subsection and subsubsection entries in the TOC
\setcounter{tocdepth}{3}
% 插入目录
\tableofcontents

\renewcommand{\theequation}{S\arabic{equation}}
\setcounter{equation}{0}

\renewcommand{\thefigure}{S\arabic{figure}}
\setcounter{figure}{0}




\newpage

% Supplementary Note 1
% Suggested packages: amsmath, amssymb, booktabs, graphicx, float
% All action and reputation labels are set upright. The symbols b and c are
% reserved exclusively for the benefit and cost of cooperation.


\section{Model formulation and evolutionary dynamics}
\label{sec:si-model}

\subsection{Donation game and behavioural strategies}
\label{sec:si-game}

We consider a large, well-mixed population. In each interaction, one individual is randomly assigned the role of donor and another the role of recipient. The donor chooses whether to cooperate ($\mathrm{C}$) or defect ($\mathrm{D}$). Cooperation gives the recipient a benefit $b$ at a cost $c$ to the donor, where $b>c>0$. Defection has no payoff consequence. With rows and columns corresponding to $\mathrm{C}$ and $\mathrm{D}$, respectively, the payoff matrix is
\[
\begin{pmatrix}
b-c & -c\\
b & 0
\end{pmatrix}.
\]

Individuals use one of three behavioural strategies. Unconditional cooperators ($\mathrm{ALLC}$) intend to cooperate with every recipient, whereas unconditional defectors ($\mathrm{ALLD}$) always intend to defect. Discriminators ($\mathrm{DISC}$) condition their intended action on the recipient's reputation. They intend to cooperate with recipients they regard as good and to defect against recipients they regard as bad.

An interaction proceeds in the following order. The donor first consults its private reputation of the recipient and chooses an intended action according to its behavioural strategy. An implementation error may alter this action. An observer then evaluates the realized action using its own private reputation of the recipient, and an assessment error may reverse this evaluation. The resulting positive or negative assessment updates the observer's internal reputation state for the donor. The following sections specify each step of this process.

\subsection{Private assessment, social norms and errors}
\label{sec:si-private-assessment}

We adopt a private-reputation framework. Each individual maintains a personal reputation record of every other member of the population. Consequently, different observers may hold different opinions of the same individual.

\textbf{Social norms.} Observers use the same assessment rule. We represent this rule in the standard second-order form, in which the donor's realized action and the observer's private reputation of the recipient jointly determine the assessment. Let $d_{A,R}$ denote the probability of a positive assessment after action $A\in\{\mathrm{C},\mathrm{D}\}$ towards a recipient whom the observer regards as $R\in\{\mathrm{G},\mathrm{B}\}$. All norms considered here assign a positive assessment to cooperation with a good recipient and a negative assessment to defection against a good recipient. They differ in how they assess actions towards a bad recipient. Let $p$ be the probability of a positive assessment after cooperation with a bad recipient, and let $q$ be the corresponding probability after defection. For the four deterministic norms considered here, $p,q\in\{0,1\}$. The four entries of the norm are therefore
\[
d_{\mathrm{C},\mathrm{G}}=1,\qquad
d_{\mathrm{D},\mathrm{G}}=0,\qquad
d_{\mathrm{C},\mathrm{B}}=p,\qquad
d_{\mathrm{D},\mathrm{B}}=q.
\]
The choices of $p$ and $q$ for the four assessment rules used in the main text are listed in Supplementary Table~\ref{tab:si-norms}.

\begin{table}[H]
\centering
\caption{\textbf{Social norms.} The parameters $p$ and $q$ are the probabilities of prescribing a positive assessment after cooperation and defection towards a bad recipient, respectively.}
\label{tab:si-norms}
\begin{tabular}{lccc}
\toprule
Norm & Abbreviation & $p$ & $q$\\
\midrule
Scoring & SC & 1 & 0\\
Shunning & SH & 0 & 0\\
Stern Judging & SJ & 0 & 1\\
Simple Standing & SS & 1 & 1\\
\bottomrule
\end{tabular}
\end{table}

\textbf{Errors.} Errors can occur at two stages. First, an \emph{implementation error} occurs with probability $e_1$, where $0\leq e_1<1/2$. Intended cooperation then fails and is implemented as defection, whereas intended defection is implemented without error. Second, after the observer evaluates the realized action according to the social norm, an \emph{assessment error} reverses the resulting assessment with probability $e_2$, where $0\leq e_2<1/2$. The two error events are independent.

Let $P^+_{A\mid R}$ denote the probability that the donor ultimately receives a positive assessment, conditional on the intended action $A\in\{\mathrm{C},\mathrm{D}\}$ and the observer's reputation of the recipient $R\in\{\mathrm{G},\mathrm{B}\}$. This probability includes both implementation and assessment errors. We obtain its four values by separating the possible events as follows.

\begin{enumerate}
\item \textbf{Intended $\mathrm{C}$ towards $\mathrm{G}$.}
\begin{enumerate}
\item With probability $1-e_1$, cooperation is successfully implemented. Cooperation with a good recipient is prescribed a positive assessment, which remains positive with probability $1-e_2$.
\item With probability $e_1$, cooperation fails and is implemented as defection. Defection against a good recipient is prescribed a negative assessment, which is reversed with probability $e_2$.
\end{enumerate}
The resulting probability of a positive assessment is
\begin{equation}
    P^+_{\mathrm{C}\mid\mathrm{G}}=(1-e_1)(1-e_2)+e_1e_2\equiv \epsilon.
\end{equation}

\item \textbf{Intended $\mathrm{D}$ towards $\mathrm{G}$.}

Intended defection is implemented correctly. Because defection against a good recipient is prescribed a negative assessment, the donor receives a positive assessment only if assessment error reverses the prescribed judgment. Hence,
\begin{equation}
    P^+_{\mathrm{D}\mid\mathrm{G}}=e_2.
\end{equation}

\item \textbf{Intended $\mathrm{C}$ towards $\mathrm{B}$.}

\begin{enumerate}
\item With probability $1-e_1$, cooperation is successfully implemented. With probability $p$, cooperation with a bad recipient is prescribed a positive assessment, which remains positive with probability $1-e_2$.

\item With probability $1-e_1$, cooperation is successfully implemented. With probability $1-p$, it is prescribed a negative assessment, which is reversed with probability $e_2$.

\item With probability $e_1$, cooperation fails and is implemented as defection. With probability $q$, defection against a bad recipient is prescribed a positive assessment, which remains positive with probability $1-e_2$.

\item With probability $e_1$, cooperation fails and is implemented as defection. With probability $1-q$, it is prescribed a negative assessment, which is reversed with probability $e_2$.
\end{enumerate}

Summing these four mutually exclusive events gives
\begin{equation}
    \begin{aligned}
    P^+_{\mathrm{C}\mid\mathrm{B}}
    ={}&(1-e_1)
    \left[p(1-e_2)+(1-p)e_2\right]+e_1
    \left[q(1-e_2)+(1-q)e_2\right] \\
    ={}&p(\epsilon-e_2)
    +q(1-\epsilon-e_2)+e_2.
    \end{aligned}
\end{equation}

\item \textbf{Intended $\mathrm{D}$ towards $\mathrm{B}$.}

\begin{enumerate}
\item With probability $q$, defection against a bad recipient is prescribed a positive assessment, which remains positive with probability $1-e_2$.

\item With probability $1-q$, it is prescribed a negative assessment, which is reversed with probability $e_2$.
\end{enumerate}

Therefore,
\begin{equation}
    P^+_{\mathrm{D}\mid\mathrm{B}}=q(1-e_2)+(1-q)e_2=q(1-2e_2)+e_2.
\end{equation}
\end{enumerate}

The corresponding probabilities for the four social norms are summarized in Supplementary Table~\ref{tab:si-effective-assessment}.

\begin{table}[H]
\centering
\caption{\textbf{Effective assessment probabilities.} Each entry is the probability of assigning a positive assessment, conditional on the donor's intended action and the observer's reputation of the recipient.}
\label{tab:si-effective-assessment}
\begin{tabular}{lcccc}
\toprule
Norm
& $P^+_{\mathrm{C}\mid\mathrm{G}}$
& $P^+_{\mathrm{D}\mid\mathrm{G}}$
& $P^+_{\mathrm{C}\mid\mathrm{B}}$
& $P^+_{\mathrm{D}\mid\mathrm{B}}$\\
\midrule
Scoring & $\epsilon$ & $e_2$ & $\epsilon$ & $e_2$\\
Shunning & $\epsilon$ & $e_2$ & $e_2$ & $e_2$\\
Stern Judging & $\epsilon$ & $e_2$ & $1-\epsilon$ & $1-e_2$\\
Simple Standing & $\epsilon$ & $e_2$ & $1-e_2$ & $1-e_2$\\
\bottomrule
\end{tabular}
\end{table}

\subsection{Inertial reputation updating}
\label{sec:si-inertia}

For each observer $i$ and target $j\neq i$, let $L_{ij}$ denote observer $i$'s internal reputation state for target $j$. Different observers maintain separate records of the same target, so $L_{ij}$ and $L_{i'j}$ need not coincide. The state space consists of two endpoint states and two paths of intermediate states,
\[
\mathcal L=\{\mathrm{G},\mathrm{W}_1,\ldots,\mathrm{W}_{k_1},
\mathrm{B},\mathrm{S}_1,\ldots,\mathrm{S}_{k_2}\}.
\]
Here $k_1,k_2\in\mathbb N$ are positive integers that specify the numbers of warning and suspicious states, respectively. The warning states form the path from established good to bad, whereas the suspicious states form the reverse path (Supplementary Fig.~\ref{fig:si-state-transitions}).

\begin{figure}[H]
\centering
\includegraphics[width=0.88\linewidth]{FigS1_inertial_state_transitions.png}
\caption{\textbf{Inertial reputation updating.} Negative assessments move an internal reputation from $\mathrm{G}$ through the warning states towards $\mathrm{B}$, whereas a positive assessment clears a warning and returns the state to $\mathrm{G}$. Positive assessments move an internal reputation from $\mathrm{B}$ through the suspicious states towards $\mathrm{G}$, whereas a negative assessment clears a suspicious state and returns it to $\mathrm{B}$. Positive assessments leave $\mathrm{G}$ unchanged, and negative assessments leave $\mathrm{B}$ unchanged.}
\label{fig:si-state-transitions}
\end{figure}

The state transitions can be summarized as follows.
\begin{itemize}
\item Following a negative assessment, $\mathrm{G}$ moves to $\mathrm{W}_1$, $\mathrm{W}_m$ moves to $\mathrm{W}_{m+1}$ for $1\leq m<k_1$, and $\mathrm{W}_{k_1}$ moves to $\mathrm{B}$. A negative assessment leaves $\mathrm{B}$ unchanged and returns any suspicious state directly to $\mathrm{B}$. A transition from $\mathrm{G}$ to $\mathrm{B}$ therefore requires $k_1+1$ consecutive negative assessments.

\item Following a positive assessment, $\mathrm{B}$ moves to $\mathrm{S}_1$, $\mathrm{S}_n$ moves to $\mathrm{S}_{n+1}$ for $1\leq n<k_2$, and $\mathrm{S}_{k_2}$ moves to $\mathrm{G}$. A positive assessment leaves $\mathrm{G}$ unchanged and returns any warning state directly to $\mathrm{G}$. A transition from $\mathrm{B}$ to $\mathrm{G}$ therefore requires $k_2+1$ consecutive positive assessments.
\end{itemize}

Behaviour depends on a binary reputation rather than directly on the internal state. Whenever the reputation is consulted, $\mathrm{G}$ is read as good and $\mathrm{B}$ as bad. A warning state $\mathrm{W}_m$ is read as good with probability $\alpha^m$ ($0\leq \alpha \leq 1$), and a suspicious state $\mathrm{S}_n$ is read as good with probability $1-\beta^n$ $0\leq \beta \leq 1$. Let $R_{ij}=1$ when observer $i$ reads target $j$ as good, we have
\begin{equation}
\Pr(R_{ij}=1\mid L_{ij}=\ell)=
\begin{cases}
1, & \ell=\mathrm{G},\\
\alpha^m, & \ell=\mathrm{W}_m,\\
0, & \ell=\mathrm{B},\\
1-\beta^n, & \ell=\mathrm{S}_n.
\end{cases}
\label{eq:si-binary-readout}
\end{equation}
The internal state persists between assessment events. Its binary reputation is sampled when it is consulted. Increasing $\alpha$ makes warning states more likely to be treated as good, whereas increasing $\beta$ makes suspicious states more likely to be treated as bad. The one-step model in the main text has $k_1=k_2=1$, with the four internal states $\mathrm{G}$, $\mathrm{W}$, $\mathrm{B}$ and $\mathrm{S}$.

We first characterize the updating rule independently of the population composition. Suppose that, while the assessment environment is held fixed, an assessment of a focal individual is positive with probability $\varphi$ and negative with probability $1-\varphi$. Successive assessments move the individual's internal state according to Supplementary Fig.~\ref{fig:si-state-transitions}, thereby defining a finite, time-homogeneous Markov chain.

We order the states as $\mathrm{G},\mathrm{W}_1,\ldots,\mathrm{W}_{k_1},\mathrm{B},\mathrm{S}_1,\ldots,\mathrm{S}_{k_2}$. Let $\mathbf M^+$ be the transition matrix conditional on a positive assessment. Its entry $[\mathbf M^+]_{\ell\ell'}$ equals one if a positive assessment moves state $\ell$ to state $\ell'$, and zero otherwise. The matrix $\mathbf M^-$ is defined in the same way for a negative assessment. Thus, each row of either conditional matrix contains exactly one nonzero entry, and the transition matrix without conditioning on the sign of the assessment is
\begin{equation}
\mathbf M(\varphi)=\varphi\mathbf M^+ +(1-\varphi)\mathbf M^-,
\label{eq:si-markov-matrix}
\end{equation}
because the two conditional transitions occur with probabilities $\varphi$ and $1-\varphi$, respectively.

Let $\boldsymbol\pi(\varphi)$ denote the stationary distribution of the Markov chain in Eq.~\eqref{eq:si-markov-matrix}. In the stationary distribution, the probabilities along the warning path satisfy $\pi_{\mathrm{W}_m}=(1-\varphi)^m\pi_{\mathrm{G}}$, and those along the suspicious path satisfy $\pi_{\mathrm{S}_n}=\varphi^n\pi_{\mathrm{B}}$. The stationary flow from the warning path into $\mathrm{B}$ must equal the flow from the suspicious path into $\mathrm{G}$. Hence,
\[
(1-\varphi)\pi_{\mathrm{W}_{k_1}}=\varphi\pi_{\mathrm{S}_{k_2}},
\qquad\text{or equivalently}\qquad
(1-\varphi)^{k_1+1}\pi_{\mathrm{G}}=\varphi^{k_2+1}\pi_{\mathrm{B}}.
\]
Together with the normalization condition
\[
\pi_{\mathrm{G}}+\sum_{m=1}^{k_1}\pi_{\mathrm{W}_m}+\pi_{\mathrm{B}}+\sum_{n=1}^{k_2}\pi_{\mathrm{S}_n}=1,
\]
these relations give
\begin{equation}
\begin{aligned}
\pi_{\mathrm{G}}(\varphi)
&=\frac{\varphi^{k_2+1}}
{(1-\varphi)^{k_1}+\varphi^{k_2}-(1-\varphi)^{k_1}\varphi^{k_2}},\\
\pi_{\mathrm{W}_m}(\varphi)
&=\frac{(1-\varphi)^m\varphi^{k_2+1}}
{(1-\varphi)^{k_1}+\varphi^{k_2}-(1-\varphi)^{k_1}\varphi^{k_2}},\\
\pi_{\mathrm{B}}(\varphi)
&=\frac{(1-\varphi)^{k_1+1}}
{(1-\varphi)^{k_1}+\varphi^{k_2}-(1-\varphi)^{k_1}\varphi^{k_2}},\\
\pi_{\mathrm{S}_n}(\varphi)
&=\frac{\varphi^n(1-\varphi)^{k_1+1}}
{(1-\varphi)^{k_1}+\varphi^{k_2}-(1-\varphi)^{k_1}\varphi^{k_2}},
\end{aligned}
\label{eq:si-stationary-distribution}
\end{equation}
where $1\leq m\leq k_1$ and $1\leq n\leq k_2$. The appearance of $k_2$ in the numerator of $\pi_{\mathrm{G}}$ and $k_1$ in the numerator of $\pi_{\mathrm{B}}$ reflects the two transition paths. Moving from $\mathrm{B}$ to $\mathrm{G}$ requires $k_2+1$ consecutive positive assessments, whereas moving from $\mathrm{G}$ to $\mathrm{B}$ requires $k_1+1$ consecutive negative assessments.

Once the Markov chain has reached its stationary distribution, the internal state is converted into a binary reputation according to Eq.~\eqref{eq:si-binary-readout}. We denote the resulting probability of a good reputation by $\mathcal R(\varphi)$. Weighting the stationary probabilities in Eq.~\eqref{eq:si-stationary-distribution} by their respective good-readout probabilities gives
\begin{equation}
\begin{aligned}
\mathcal R(\varphi)
&=\pi_{\mathrm{G}}(\varphi)+\sum_{m=1}^{k_1}\pi_{\mathrm{W}_m}(\varphi)\alpha^m +\sum_{n=1}^{k_2}\pi_{\mathrm{S}_n}(\varphi)(1-\beta^n)\\
&=
\frac{\varphi^{k_2+1}\left[1+\displaystyle\sum_{m=1}^{k_1}\{\alpha(1-\varphi)\}^m\right]+(1-\varphi)^{k_1+1}\displaystyle\sum_{n=1}^{k_2}\varphi^n(1-\beta^n)}{(1-\varphi)^{k_1}+\varphi^{k_2}-(1-\varphi)^{k_1}\varphi^{k_2}}.
\end{aligned}
\label{eq:si-reputation-transform}
\end{equation}
For the one-step model, the difference between the equilibrium good reputation and the positive-assessment probability simplifies to
\begin{equation}
\mathcal R(\varphi)-\varphi
=\frac{\varphi(1-\varphi)}{1-\varphi+\varphi^2}
\left[\alpha\varphi-\beta(1-\varphi)\right].
\label{eq:si-one-step-correction}
\end{equation}
Tolerance therefore raises the probability of a good reputation through the term $\alpha\varphi$, whereas strictness lowers it through the term $\beta(1-\varphi)$. The baseline identity extends beyond the one-step model. When $\alpha=\beta=0$, every warning state is read as bad and every suspicious state as good. A positive assessment then produces a good binary reputation, whereas a negative assessment produces a bad binary reputation, irrespective of the preceding internal state or the values of $k_1$ and $k_2$. The latest assessment therefore directly determines the binary reputation, as under traditional instantaneous updating, and $\mathcal R(\varphi)=\varphi$ for all non-negative integers $k_1$ and $k_2$. For $k_1=k_2=1$, the same result follows immediately from Eq.~\eqref{eq:si-one-step-correction}.

% Tolerance therefore raises the probability of a good reputation through the term $\alpha\varphi$, whereas strictness lowers it through the term $\beta(1-\varphi)$. For the one-step model with $\alpha=\beta=0$, $\mathrm{W}$ is always read as bad and $\mathrm{S}$ is always read as good. A negative assessment therefore produces a bad binary reputation, and a positive assessment produces a good binary reputation, irrespective of the preceding internal state. The latest assessment directly replaces the binary reputation, as under traditional instantaneous updating, and Eq.~\eqref{eq:si-one-step-correction} reduces to $\mathcal R(\varphi)=\varphi$.



\subsection{Full-state reputation dynamics and stationary reduction}
\label{sec:si-reputation-dynamics}
\label{sec:si-full-state-dynamics}

We next determine the reputations generated by a population with a fixed strategy composition. Let $\mathcal S=\{\mathrm{ALLC},\mathrm{ALLD},\mathrm{DISC}\}$ and let $f_I$ denote the frequency of strategy $I\in\mathcal S$, where $\sum_{I\in\mathcal S}f_I=1$. We write $r_I$ for the probability that a randomly chosen observer regards a type-$I$ individual as good. The probability that a randomly chosen recipient is regarded as good is therefore $\bar r=\sum_{I\in\mathcal S}f_Ir_I$.

A $\mathrm{DISC}$ donor acts according to its own reputation of the recipient, whereas the donor is assessed according to the observer's reputation of that recipient. Following standard mean-field treatments of private assessment, we assume that these two reputations are conditionally independent given the recipient's strategy. The probabilities that both individuals regard a random recipient as good, both regard the recipient as bad, or hold a particular ordered pair of disagreeing opinions are, respectively,
\begin{equation}
\begin{aligned}
h_{\mathrm{GG}}&=\sum_{I\in\mathcal S}f_Ir_I^2,\\
h_{\mathrm{BB}}&=\sum_{I\in\mathcal S}f_I(1-r_I)^2,\\
h_{\mathrm{GB}}&=\sum_{I\in\mathcal S}f_Ir_I(1-r_I).
\end{aligned}
\label{eq:si-private-pairs}
\end{equation}
Here $h_{\mathrm{GB}}$ applies to either ordered disagreement, so the total probability that the donor and observer disagree is $2h_{\mathrm{GB}}$.

Let $\varphi_I$ denote the probability that a type-$I$ donor receives a positive assessment. These probabilities follow from the action rule of each strategy and the four effective assessment probabilities in Supplementary Table~\ref{tab:si-effective-assessment}.
\begin{itemize}
\item An $\mathrm{ALLC}$ donor intends to cooperate regardless of the recipient's reputation. The observer regards the recipient as good with probability $\bar r$ and as bad with probability $1-\bar r$. Thus, the two possible contributions to a positive assessment are $\bar rP^+_{\mathrm{C}\mid\mathrm{G}}$ and $(1-\bar r)P^+_{\mathrm{C}\mid\mathrm{B}}$.

\item An $\mathrm{ALLD}$ donor always intends to defect. The corresponding contributions are $\bar rP^+_{\mathrm{D}\mid\mathrm{G}}$ and $(1-\bar r)P^+_{\mathrm{D}\mid\mathrm{B}}$.

\item A $\mathrm{DISC}$ donor intends to cooperate if it regards the recipient as good and to defect otherwise. There are four cases. If donor and observer both regard the recipient as good, which occurs with probability $h_{\mathrm{GG}}$, the relevant assessment probability is $P^+_{\mathrm{C}\mid\mathrm{G}}$. If the donor regards the recipient as good but the observer regards the recipient as bad, the contribution is $h_{\mathrm{GB}}P^+_{\mathrm{C}\mid\mathrm{B}}$. If the donor regards the recipient as bad but the observer regards the recipient as good, the contribution is $h_{\mathrm{GB}}P^+_{\mathrm{D}\mid\mathrm{G}}$. Finally, if both regard the recipient as bad, the contribution is $h_{\mathrm{BB}}P^+_{\mathrm{D}\mid\mathrm{B}}$.
\end{itemize}

Combining these mutually exclusive cases gives
\begin{equation}
\begin{aligned}
\varphi_{\mathrm{ALLC}}
&=\bar rP^+_{\mathrm{C}\mid\mathrm{G}}
+(1-\bar r)P^+_{\mathrm{C}\mid\mathrm{B}},\\
\varphi_{\mathrm{ALLD}}
&=\bar rP^+_{\mathrm{D}\mid\mathrm{G}}
+(1-\bar r)P^+_{\mathrm{D}\mid\mathrm{B}},\\
\varphi_{\mathrm{DISC}}
&=h_{\mathrm{GG}}P^+_{\mathrm{C}\mid\mathrm{G}}
+h_{\mathrm{GB}}\left(P^+_{\mathrm{C}\mid\mathrm{B}}
+P^+_{\mathrm{D}\mid\mathrm{G}}\right)
+h_{\mathrm{BB}}P^+_{\mathrm{D}\mid\mathrm{B}}.
\end{aligned}
\label{eq:si-positive-assessment-probabilities}
\end{equation}

For each target strategy $I$, let $\boldsymbol\pi_I(t)=(\pi_{I,\mathrm G},\pi_{I,\mathrm W_1},\ldots,\pi_{I,\mathrm W_{k_1}},\pi_{I,\mathrm B},\pi_{I,\mathrm S_1},\ldots,\pi_{I,\mathrm S_{k_2}})$ be the row vector of internal-state probabilities. With $d=2+k_1+k_2$, $\boldsymbol\pi_I$ has dimension $1\times d$, the transition matrix $\mathbf M$ has dimension $d\times d$, and the readout vector below has dimension $d\times1$:
\begin{equation}
\boldsymbol\omega=(1,\alpha,\ldots,\alpha^{k_1},0,1-\beta,\ldots,1-\beta^{k_2})^{\mathsf T},
\qquad r_I(t)=\boldsymbol\pi_I(t)\boldsymbol\omega.
\label{eq:si-full-state-readout}
\end{equation}
A zero-length path contributes no intermediate entries. The assessment probabilities $\varphi_I(\boldsymbol r;\boldsymbol f)$ are those in Eq.~\eqref{eq:si-positive-assessment-probabilities}.

The simulations update one observer--target record at a time. The corresponding asynchronous mean-field description uses a continuous time scaled so that a record is updated at unit rate:
\begin{equation}
\frac{\mathrm d\boldsymbol\pi_I}{\mathrm dt}
=\boldsymbol\pi_I\{\mathbf M[\varphi_I(\boldsymbol r;\boldsymbol f)]-\mathbf I\}.
\label{eq:si-full-state-master}
\end{equation}
This equation preserves non-negativity and unit total probability. It retains the same conditional-independence closure as the main-text reputation equations; it is not an exact equation for the joint private-image distribution of a finite population. At a finite-population event, the conditional expected increment of the empirical state distribution has the corresponding factor $1/[n_{\mathrm{DISC}}N]$, with assessment probabilities depending on the full private-image configuration. Replacing those probabilities by the mean-field closure and measuring time in sweeps gives Eq.~\eqref{eq:si-full-state-master}.

Because $\boldsymbol\omega$ is fixed, projecting Eq.~\eqref{eq:si-full-state-master} gives the exact mean-field readout dynamics
\begin{equation}
\frac{\mathrm d r_I}{\mathrm dt}
=\frac{\mathrm d\boldsymbol\pi_I}{\mathrm dt}\boldsymbol\omega
=\boldsymbol\pi_I\mathbf M[\varphi_I(\boldsymbol r;\boldsymbol f)]\boldsymbol\omega-r_I.
\label{eq:si-full-state-projection}
\end{equation}
Although $\varphi_I$ depends on $\boldsymbol r$, this dependence introduces no additional chain-rule term in Eq.~\eqref{eq:si-full-state-projection}: the differentiated definition is $r_I=\boldsymbol\pi_I\boldsymbol\omega$, not a composition involving $\mathbf M$. The current value of $\varphi_I(\boldsymbol r;\boldsymbol f)$ is simply evaluated in the vector field. Derivatives of $\mathbf M[\varphi_I(\boldsymbol r;\boldsymbol f)]$ enter only when differentiating that vector field, for example when constructing its Jacobian for a stability calculation.

The first term in Eq.~\eqref{eq:si-full-state-projection} is the expected readout after one update from the current latent-state distribution. It is generally not the stationary response $\mathcal R(\varphi_I)$. A direct counterexample is the one-step model with $k_1=k_2=1$ and $\alpha=\beta=1$. Starting from $\mathrm G$, one update leads to either $\mathrm G$ after a positive assessment or $\mathrm W$ after a negative assessment, and both states are read as good. Therefore, for $\boldsymbol\pi_{\mathrm G}=(1,0,0,0)$,
\begin{equation}
\boldsymbol\pi_{\mathrm G}\mathbf M(u)\boldsymbol\omega=1,
\qquad
\mathcal R(u)=\frac{u^2(2-u)}{1-u+u^2}<1
\quad (0<u<1).
\label{eq:si-one-step-versus-stationary}
\end{equation}
Thus $\boldsymbol\pi_I\mathbf M(\varphi_I)\boldsymbol\omega$ retains information about the current distribution over $\mathrm G$, warning, $\mathrm B$ and suspicious states, whereas $\mathcal R(\varphi_I)$ first replaces that distribution by the stationary distribution for a fixed assessment probability.

For fixed exogenous $u$, the embedded chain of a focal record at successive observation times is exactly $\boldsymbol\pi_{m+1}=\boldsymbol\pi_m\mathbf M(u)$. Substituting an endogenous $\varphi_I(\boldsymbol r_m)$ and updating all distributions simultaneously instead defines a synchronous mean-field map. It has the same stationary equations, but its transient dynamics and stability need not coincide with asynchronous updating. We therefore use Eq.~\eqref{eq:si-full-state-master} for the full-state branch check.

At stationarity,
\begin{equation}
\boldsymbol\pi_I^*\mathbf M(\varphi_I^*)=\boldsymbol\pi_I^*,
\qquad \boldsymbol\pi_I^*\boldsymbol 1=1.
\label{eq:si-full-state-stationary}
\end{equation}
For $0<\varphi_I^*<1$, the direct-reset chain is irreducible and aperiodic. Its unique stationary distribution is Eq.~\eqref{eq:si-stationary-distribution}; the endpoints have the corresponding unique absorbing-state limits. Consequently,
\begin{equation}
r_I^*=\boldsymbol\pi_I^*\boldsymbol\omega=\mathcal R(\varphi_I^*).
\label{eq:si-full-state-consistency}
\end{equation}
Combining this identity with the assessment closure $\varphi_I^*=\varphi_I(\boldsymbol r^*;\boldsymbol f)$ gives the exact stationary self-consistency equations
\begin{equation}
r_I^*=\mathcal R[\varphi_I(\boldsymbol r^*;\boldsymbol f)],
\qquad I\in\mathcal S.
\label{eq:si-reputation-fixed-point}
\end{equation}
Conversely, any solution of Eq.~\eqref{eq:si-reputation-fixed-point} constructs a full-state stationary solution by setting $\boldsymbol\pi_I^*=\boldsymbol\pi[\varphi_I(\boldsymbol r^*;\boldsymbol f)]$. Hence the full-state and reduced descriptions have exactly the same stationary equations within the same closure.

For low-dimensional solution and branch selection, we use the auxiliary fixed-point relaxation
\begin{equation}
\frac{\mathrm dr_I}{\mathrm d\tau}
=\mathcal R[\varphi_I(\boldsymbol r;\boldsymbol f)]-r_I,
\qquad I\in\mathcal S.
\label{eq:si-reputation-dynamics}
\end{equation}
Here $\tau$ is an auxiliary variable, not the observation time $t$. Its fixed points are exactly Eq.~\eqref{eq:si-reputation-fixed-point}, but its transient trajectories, local stability spectra and attraction basins need not match those of Eq.~\eqref{eq:si-full-state-master}. We assume that all individuals are initially regarded as good, so the auxiliary calculation starts from $r_I(0)=1$ and the full-state calculation starts from $\boldsymbol\pi_I(0)=(1,0,\ldots,0)$. Under the separation-of-timescales assumption used below, only the stationary reputations are passed to the slower strategy dynamics, for which Eq.~\eqref{eq:si-reputation-dynamics} is therefore sufficient after its selected branch has been checked against the full-state process. Temporal, periodically forced or history-dependent analyses must instead evolve the latent-state process itself.

For a pure discriminator population, $\varphi_{\mathrm{DISC}}=\Phi_R(r)$ and the condition reduces to $r^*=\mathcal R[\Phi_R(r^*)]$. We compare direct fixed-point solutions, the auxiliary relaxation and the full-state equation from their respective all-good initial conditions. The validation tables report both dynamical residuals, fixed-point residuals, convergence flags and differences between the selected branches. A direct root solver verifies a root but does not itself define a dynamically selected branch. Agreement at tested parameters is a numerical consistency check, not a general theorem of branch equivalence.

The completed audit contains 720 cases: representative readout endpoints and error combinations, together with every distinct depth pair in the existing depth tables tested in pure DISC and equal-frequency mixed populations. All cases met the numerical convergence checks. The maximum difference between auxiliary and full-state equilibrium reputation vectors was $2.37\times10^{-8}$, and the largest reported stationary residual was $1.29\times10^{-10}$. These tests support retaining the existing stationary framework for the audited cases; they do not establish branch equivalence at every point of a dense parameter scan.
\subsection{Payoffs and replicator dynamics}
\label{sec:si-evolution}

Following previous work on indirect reciprocity, we assume a separation of timescales such that reputations reach equilibrium before strategy evolution. For each strategy composition, we therefore first determine the stable solution of Eq.~\eqref{eq:si-reputation-dynamics} and then evaluate evolutionary payoffs. Here $r_I^*$ is the probability that a type-$I$ recipient is regarded as good, and $\bar r^*=\sum_{I\in\mathcal S}f_Ir_I^*$ is the corresponding probability for a randomly selected recipient.

The payoff of a strategy can be separated into incoming and outgoing cooperation. A type-$I$ recipient receives intended cooperation with probability
\[
\kappa_I^{\mathrm{in}}=f_{\mathrm{ALLC}}+f_{\mathrm{DISC}}r_I^*,
\]
because $\mathrm{ALLC}$ donors always intend to cooperate and $\mathrm{DISC}$ donors do so when they regard the recipient as good. The probabilities that donors using the three strategies intend to cooperate are
\[
\kappa_{\mathrm{ALLC}}^{\mathrm{out}}=1,
\qquad
\kappa_{\mathrm{ALLD}}^{\mathrm{out}}=0,
\qquad
\kappa_{\mathrm{DISC}}^{\mathrm{out}}=\bar r^*.
\]
Because intended cooperation is implemented with probability $1-e_1$, the payoff of strategy $I$ can be written compactly as $\Pi_I=(1-e_1)(b\kappa_I^{\mathrm{in}}-c\kappa_I^{\mathrm{out}})$. Substituting the incoming and outgoing cooperation probabilities gives
\begin{equation}
\begin{aligned}
\Pi_{\mathrm{ALLC}}
&=(1-e_1)\left[b\left(f_{\mathrm{ALLC}}+f_{\mathrm{DISC}}r_{\mathrm{ALLC}}^*\right)-c\right],\\
\Pi_{\mathrm{ALLD}}
&=(1-e_1)b\left(f_{\mathrm{ALLC}}+f_{\mathrm{DISC}}r_{\mathrm{ALLD}}^*\right),\\
\Pi_{\mathrm{DISC}}
&=(1-e_1)\left[b\left(f_{\mathrm{ALLC}}+f_{\mathrm{DISC}}r_{\mathrm{DISC}}^*\right)-c\bar r^*\right].
\end{aligned}
\label{eq:si-payoffs}
\end{equation}

On the slower evolutionary timescale, strategy frequencies follow the replicator equations
\begin{equation}
\frac{\mathrm df_I}{\mathrm dt}
=f_I(\Pi_I-\bar\Pi),
\qquad
\bar\Pi=\sum_{J\in\mathcal S}f_J\Pi_J.
\label{eq:si-replicator-dynamics}
\end{equation}
At each strategy composition, the stable reputation state is recomputed before the payoffs are evaluated. The replicator dynamics therefore describe strategy evolution under the reputation environment generated by the current population.

\subsubsection{Numerical integration}

The final Python implementation uses adaptive DOP853 integration for both timescales. The auxiliary reputation system is initialized from the all-good branch and uses \texttt{rtol}$=10^{-9}$, \texttt{atol}$=10^{-11}$, a maximum auxiliary time of $2\times10^4$ and the convergence condition $\|\dot{\boldsymbol r}\|_\infty\leq10^{-10}$. The replicator system uses \texttt{rtol}$=10^{-8}$, \texttt{atol}$=10^{-10}$, maximum step 20, maximum time $2\times10^4$ and $\|\dot{\boldsymbol f}\|_\infty\leq10^{-7}$. For the basin calculations, the stationary strategy field is evaluated on a triangular lattice, interpolated piecewise linearly, and checked by recomputing the stationary reputations at terminal compositions. Every parameter set uses the same 512 scrambled-Sobol initial compositions.

The previous publication scripts used fourth-order Runge--Kutta steps of 0.1 for the auxiliary reputation equation and midpoint Runge--Kutta steps of 0.1 for the interpolated strategy field. In 16 representative reputation cases, adaptive DOP853 and the previous implementation differed by at most $6.53\times10^{-9}$, while DOP853 and Radau differed by at most $6.59\times10^{-10}$. Across four representative strategy cases, fixed-step RK2 and DOP853 trajectories evaluated at the same time horizon differed by at most $4.33\times10^{-5}$; after convergence, DOP853 and Radau differed by at most $2.59\times10^{-10}$. All adaptive test cases met the stated residual criteria. These checks support DOP853 as the final batched Python protocol while retaining Radau as an independent regression solver.


%————————————————————————————————————————————————————————————————————
%————————————————————————————————————————————————————————————————————

\section{Evolutionary stability and resistance to invasion}
\label{sec:si-evolutionary-stability}

We next ask whether an established population of discriminators (DISC) can resist rare behavioural mutants.
Rather than treating different mutant strategies separately, we first write their invasion condition in a common form.
This allows unconditional mutants and temporally structured opportunists to be analysed within the same resident reputation environment.
Throughout this section, resident reputations are selected by relaxation from the all-good initial condition.

\subsection{Resident DISC environment and unified invasion condition}

Consider a resident population consisting entirely of DISC.
For a candidate resident reputation $u$, the probability of a positive assessment is
\begin{equation}
\Phi_{\mathrm R}(u)=u^2P^+_{\mathrm C\mid\mathrm G}+u(1-u)\left(P^+_{\mathrm C\mid\mathrm B}+P^+_{\mathrm D\mid\mathrm G}\right)+(1-u)^2P^+_{\mathrm D\mid\mathrm B}.
\label{eq:si-resident-assessment-strategy}
\end{equation}
The selected resident reputation $r_R^*$ satisfies
\begin{equation}
r_{\mathrm R}^{*}=\mathcal R\!\left[\Phi_{\mathrm R}(r_{\mathrm R}^{*})\right],
\label{eq:si-resident-reputation-strategy}
\end{equation}
where $\mathcal R(\phi)$ is the reputation map derived in Supplementary Note 1.



In this resident environment, an individual that always intends cooperation or defection receives a positive assessment with probabilities
\begin{align}
\Phi_{\mathrm C}^{*}
&=r_{\mathrm R}^{*}P^+_{\mathrm C\mid\mathrm G}+(1-r_{\mathrm R}^{*})P^+_{\mathrm C\mid\mathrm B},
\label{eq:si-allc-assessment-strategy}
\\
\Phi_{\mathrm D}^{*}
&=r_{\mathrm R}^{*}P^+_{\mathrm D\mid\mathrm G}+(1-r_{\mathrm R}^{*})P^+_{\mathrm D\mid\mathrm B}.
\label{eq:si-alld-assessment-strategy}
\end{align}

The resident assessment probability can equivalently be written as
\begin{equation}
\Phi_\mathrm R^*
\equiv
\Phi_\mathrm R(r_\mathrm R^*)=r_\mathrm R^*\Phi_\mathrm C^*+(1-r_\mathrm R^*)\Phi_\mathrm D^*.
\end{equation}
We define the assessment contrast
\begin{equation}
\Delta\Phi^*=\Phi_\mathrm C^*-\Phi_\mathrm D^*.
\label{eq:SI_assessment_contrast}
\end{equation}
For the four social norms considered here,
\begin{equation}
\Delta\Phi^*
=(\varepsilon-e_2)
\begin{cases}
1, & \text{Scoring},\\
r_\mathrm R^*, & \text{Shunning},\\
2r_\mathrm R^*-1, & \text{Stern Judging},\\
r_\mathrm R^*, & \text{Simple Standing}.
\end{cases}
\label{eq:SI_assessment_contrast_norms}
\end{equation}

Now consider an arbitrary rare mutant $\mathrm X$.
Let $r_\mathrm X$ denote its long-run probability of being read as good, and let $q_\mathrm X$ denote its long-run probability of intending cooperation. Because the mutant is rare, resident DISC donors cooperate with it with probability $r_\mathrm X$, whereas the mutant itself pays the cooperation cost with probability $q_\mathrm X$.
Its payoff difference from the resident is therefore
\begin{equation}
\frac{\lambda_\mathrm X}{(1-e_1)c}=\frac{b}{c}(r_\mathrm X-r_\mathrm R^*)+(r_\mathrm R^*-q_\mathrm X).
\label{eq:SI_general_invasion}
\end{equation}

It is useful to separate the two terms in Eq.~\eqref{eq:SI_general_invasion}.
We define
\begin{equation}
S_\mathrm X=r_\mathrm R^*-q_\mathrm X,
\qquad
D_\mathrm X=r_\mathrm R^*-r_\mathrm X,
\label{eq:SI_saved_cost_penalty}
\end{equation}
where $S_\mathrm X$ is the cooperation cost saved relative to resident DISC and $D_\mathrm X$ is the associated reputational penalty.
Equation~\eqref{eq:SI_general_invasion} then becomes
\begin{equation}
\frac{\lambda_\mathrm X}{(1-e_1)c}=S_\mathrm X-\frac{b}{c}D_\mathrm X.
\label{eq:SI_master_invasion}
\end{equation}

Thus, invasion is determined by a simple balance. A mutant benefits by reducing cooperation, but this advantage can be offset if the behavioural deviation produces a sufficiently large loss of reputation. For a cost-saving mutant with $S_\mathrm X>0$ and $D_\mathrm X>0$, define
\begin{equation}
R_\mathrm X=\frac{S_\mathrm X}{D_\mathrm X}.
\label{eq:SI_general_threshold}
\end{equation}
The mutant is then excluded when
\begin{equation}
\frac{b}{c}>R_\mathrm X.
\label{eq:SI_general_lower_condition}
\end{equation}

This form will be used below for both unconditional defection and temporally structured exploitation.

\subsection{Unconditional mutants}

The unconditional strategies follow directly from the general invasion condition.

For a rare ALLD mutant,
\begin{equation}
q_\mathrm D=0,
\qquad
r_\mathrm D^*=R(\Phi_\mathrm D^*).
\end{equation}
Hence
\begin{equation}
S_\mathrm D=r_\mathrm R^*,
\qquad
D_\mathrm D=r_\mathrm R^*-r_\mathrm D^*,
\end{equation}
and Eq.~\eqref{eq:SI_general_threshold} gives
\begin{equation}
R_\mathrm L
=
\frac{r_\mathrm R^*}{r_\mathrm R^*-r_\mathrm D^*}.
\label{eq:SI_RL}
\end{equation}
When $r_\mathrm R^*>r_\mathrm D^*$, rare ALLD is excluded if
\begin{equation}
\frac{b}{c}>R_\mathrm L.
\label{eq:SI_ALLD_condition}
\end{equation}
If $r_\mathrm R^*\le r_\mathrm D^*$, no finite $b/c$ can exclude ALLD.

For a rare ALLC mutant,
\begin{equation}
q_\mathrm C=1,
\qquad
r_\mathrm C^*=R(\Phi_\mathrm C^*).
\end{equation}
Its payoff difference can be written as
\begin{equation}
\frac{\lambda_{\mathrm{ALLC}}}{(1-e_1)c}
=
\frac{b}{c}(r_\mathrm C^*-r_\mathrm R^*)
-
(1-r_\mathrm R^*).
\end{equation}
When $r_\mathrm C^*>r_\mathrm R^*$, exclusion of ALLC therefore requires
\begin{equation}
\frac{b}{c}
<
R_\mathrm U,
\qquad
R_\mathrm U
=
\frac{1-r_\mathrm R^*}{r_\mathrm C^*-r_\mathrm R^*}.
\label{eq:SI_RU}
\end{equation}
If $r_\mathrm C^*\le r_\mathrm R^*$, ALLC imposes no finite upper threshold, apart from neutral limiting cases.

The familiar stability interval for pure DISC therefore follows directly from the unified invasion condition:
\begin{equation}
\boxed{
R_\mathrm L
<
\frac{b}{c}
<
R_\mathrm U
}.
\label{eq:SI_DISC_window}
\end{equation}
The lower boundary reflects resistance to strategies that save cooperation costs, whereas the upper boundary reflects resistance to strategies that cooperate more than resident DISC.

\subsection{Independent partial under-cooperation}
\label{sec:si-iid-opportunism}

An under-cooperation mutant follows the resident DISC rule except that, whenever its private readout of the recipient prescribes cooperation, it independently intends defection with probability $\delta\in[0,1]$. This extra decision occurs before the implementation error. It is not an unconditional probabilistic cooperator: it never adds cooperation towards a recipient it reads as bad. The mean probability of intending cooperation in the pure-DISC resident environment is
\begin{equation}
q_{\mathrm X}(\delta)=(1-\delta)r_{\mathrm R}^*.
\label{eq:si-iid-intended-cooperation}
\end{equation}
Conditional on the donor's intended action, $\Phi_{\mathrm C}^*$ and $\Phi_{\mathrm D}^*$ in Eqs.~\eqref{eq:si-allc-assessment-strategy} and \eqref{eq:si-alld-assessment-strategy} already include implementation and assessment errors. Therefore the mutant's positive-assessment probability is
\begin{align}
\Phi_{\mathrm X}(\delta)
&=(1-\delta)r_{\mathrm R}^*\Phi_{\mathrm C}^*
 +[1-(1-\delta)r_{\mathrm R}^*]\Phi_{\mathrm D}^*\\
&=\Phi_{\mathrm R}^*-\delta r_{\mathrm R}^*\Delta\Phi^*.
\label{eq:si-iid-assessment}
\end{align}
Its independent actions make successive assessments iid within the fixed rare-mutant environment, so its stationary good-reputation probability is
\begin{equation}
r_{\mathrm X}^*(\delta)=\mathcal R[\Phi_{\mathrm X}(\delta)].
\label{eq:si-iid-reputation}
\end{equation}
Substitution in the general payoff identity, Eq.~\eqref{eq:SI_master_invasion}, yields
\begin{equation}
\Lambda_{\mathrm{iid}}(\delta)
\equiv\frac{\lambda_{\mathrm X}(\delta)}{(1-e_1)c}
=\delta r_{\mathrm R}^*-\frac bc
\left\{r_{\mathrm R}^*-\mathcal R[\Phi_{\mathrm R}^*-\delta r_{\mathrm R}^*\Delta\Phi^*]\right\}.
\label{eq:si-iid-invasion}
\end{equation}
The term in braces is the reputation penalty and the first term is the cost saving. At $\delta=0$, Eq.~\eqref{eq:si-iid-invasion} is exactly zero. At $\delta=1$, $\Phi_{\mathrm X}=\Phi_{\mathrm D}^*$, so $\Lambda_{\mathrm{iid}}(1)=\lambda_{\mathrm{ALLD}}/[(1-e_1)c]$ and the classical ALLD criterion is recovered. In between, the sign is determined by the full curvature of $\mathcal R$ and need not be implied by the ALLD endpoint.

Differentiating at the resident gives the local condition
\begin{equation}
\Lambda_{\mathrm{iid}}'(0)
=r_{\mathrm R}^*\left[1-\frac bc\Delta\Phi^*\mathcal R'(\Phi_{\mathrm R}^*)\right].
\label{eq:si-iid-local-slope}
\end{equation}
A positive derivative proves invasion for sufficiently small $\delta>0$, whereas a negative derivative excludes sufficiently small deviations. At zero derivative, higher-order terms are generally needed. A non-positive derivative alone does not prove exclusion for finite $\delta$. We define complete exclusion of this \emph{tested family} by $\sup_{0<\delta\leq1}\Lambda_{\mathrm{iid}}(\delta)\leq0$, with strict negativity away from the neutral resident endpoint where appropriate. The width $\delta_{\mathrm c}$ in Extended Data Fig.~\ref{fig:ed-alpha-invasion} is the total length of intervals on which $\Lambda_{\mathrm{iid}}>0$; for the displayed representative curves, each nonempty set is a single interval beginning at zero. Readout security and the distinct high-reputation branch limit are derived next.

\subsection{Readout security and high-reputation branch limits}
\label{sec:si-readout-security}

\subsubsection{Complete-family security and reputation curvature}
\label{sec:si-readout-global}

All results below use the same private-assessment closure and rare-mutant environment as Eqs.~\eqref{eq:si-iid-assessment}--\eqref{eq:si-iid-invasion}. We write the dependence of the stationary mapping on readout explicitly as $\mathcal R_{\alpha,\beta}$, with fixed path depths $k_1,k_2$. For each parameter setting, the selected resident has reputation $r^*_{\alpha,\beta}>0$, assessment probability $u^*_{\alpha,\beta}=\Phi_{\mathrm R}(r^*_{\alpha,\beta})$ and assessment contrast $\Delta\Phi^*_{\alpha,\beta}$. These quantities must be recomputed when the readout changes. With $g_{\alpha,\beta}$ defined in Eq.~\eqref{eq:methods-readout-penalty},
\begin{equation}
\Lambda_{\mathrm{iid}}(\delta;\alpha,\beta)
=\delta r^*_{\alpha,\beta}
\left[1-\frac bc g_{\alpha,\beta}(\delta)\right].
\label{eq:si-readout-gain-identity}
\end{equation}
Extending $g$ continuously to $\delta=0$ gives the exact complete-family condition
\begin{equation}
G_{\mathrm{global}}(\alpha,\beta;b/c)
\equiv\frac bc\min_{0\leq\delta\leq1}g_{\alpha,\beta}(\delta)\geq1.
\label{eq:si-readout-global-condition}
\end{equation}
In comparison, the first-order local condition uses only
\begin{equation}
G_{\mathrm{loc}}
=\frac bc\Delta\Phi^*_{\alpha,\beta}
\mathcal R'_{\alpha,\beta}(u^*_{\alpha,\beta}).
\label{eq:si-readout-local-condition}
\end{equation}
For $G_{\mathrm{loc}}=1$, a local equality does not generally resolve either higher-order local invasion or finite deviations. The minimum in Eq.~\eqref{eq:si-readout-global-condition} retains both possibilities.

Let $\mathcal A_{\mathrm{sec}}$ be the set of admissible $\alpha\in[0,1]$ on the specified resident branch satisfying Eq.~\eqref{eq:si-readout-global-condition}. We define
\begin{equation}
\alpha_{\mathrm{sec}}=\sup\mathcal A_{\mathrm{sec}},
\label{eq:si-alpha-sec-definition}
\end{equation}
when this set is nonempty. This is the largest secure warning readout when the upper endpoint is included. The definition does not assume that every smaller $\alpha$ is secure or even belongs to the same resident branch. When $g_{\min}=\min_\delta g_{\alpha,\beta}(\delta)>0$, the corresponding critical payoff ratio is
\begin{equation}
\left(\frac bc\right)_{\mathrm{crit}}(\alpha,\beta)=\frac{1}{g_{\min}}.
\label{eq:si-readout-critical-bc}
\end{equation}
If $g_{\min}\leq0$, no finite positive payoff ratio excludes the whole family. Because the resident environment is independent of $b/c$, security sets on a fixed branch are nested as $b/c$ increases. This explains why the upper boundaries in Extended Data Fig.~\ref{fig:ed-alpha-invasion} increase with the payoff ratio without moving the resident branch limit.

A structural equivalence between local and complete-family security holds when $k_2=0$. Set $k=k_1\geq1$. Equation~\eqref{eq:si-reputation-transform} reduces to
\begin{equation}
\begin{aligned}
\mathcal R_\alpha(u)
&=u\sum_{m=0}^{k}[\alpha(1-u)]^m\\
&=(1-\alpha)\sum_{j=0}^{k-1}\alpha^j
\left[1-(1-u)^{j+1}\right]
+\alpha^k\left[1-(1-u)^{k+1}\right].
\end{aligned}
\label{eq:si-readout-concave-mixture}
\end{equation}
The weights in the second expression are nonnegative and sum to one. Each function $1-(1-u)^m$, $m\geq1$, is concave on $[0,1]$; hence $\mathcal R_\alpha''(u)\leq0$. For $\Delta\Phi^*_{\alpha,\beta}>0$, the supporting-tangent inequality implies
\begin{equation}
\mathcal R_\alpha(u^*_{\alpha,\beta}-x)
\leq\mathcal R_\alpha(u^*_{\alpha,\beta})
-x\mathcal R'_\alpha(u^*_{\alpha,\beta}),
\qquad x\geq0,
\label{eq:si-readout-tangent}
\end{equation}
whenever the argument remains in $[0,1]$. Taking $x=\delta r^*_{\alpha,\beta}\Delta\Phi^*_{\alpha,\beta}$ gives $g_{\alpha,\beta}(\delta)\geq g_{\alpha,\beta}(0)$ and therefore
\begin{equation}
G_{\mathrm{global}}=G_{\mathrm{loc}},\qquad k_2=0,quad \Delta\Phi^*_{\alpha,\beta}>0.
\label{eq:si-readout-local-global-equivalence}
\end{equation}
This proves the equivalence for the representative Scoring $(2,0)$ and Shunning $(3,0)$ configurations, including the local equality boundary. The same proof does not apply to arbitrary recovery paths. For Stern Judging $(4,3)$ and Simple Standing $(1,2)$, the full mutant interval was evaluated separately; the displayed boundaries agreed numerically with the local boundary where an interior equality existed.

\subsubsection{An exact warning-readout threshold in the error-free limit}
\label{sec:si-readout-errorfree}

Set $e_1=e_2=0$ and condition on an all-good pure-DISC resident. For all four norms, cooperation towards a good recipient is assessed positively and defection negatively, giving $r_{\mathrm R}^*=u_{\mathrm R}^*=\Delta\Phi^*=1$. Assume $\beta=1$, $k=k_1\geq1$ and $\ell=k_2\geq0$. The stationary mapping is
\begin{equation}
\mathcal R_{k,\ell,\alpha}(u)=
\frac{u^{\ell+1}\displaystyle\sum_{m=0}^{k}[\alpha(1-u)]^m}
{(1-u)^k+u^\ell-(1-u)^ku^\ell}.
\label{eq:si-readout-errorfree-map}
\end{equation}
For $v=1-u\to0$, the numerator is $1-(\ell+1-\alpha)v+O(v^2)$ and the denominator is $1-\ell v+O(v^2)$. Thus
\begin{equation}
\mathcal R_{k,\ell,\alpha}(1-v)
=1-(1-\alpha)v+O(v^2),\qquad
\mathcal R'_{k,\ell,\alpha}(1)=1-\alpha.
\label{eq:si-readout-first-warning}
\end{equation}
The first warning determines the leading response; additional warning depth affects higher-order terms. The invasion derivative is $\Lambda'_{\mathrm{iid}}(0)=1-(b/c)(1-\alpha)$, so $\alpha>1-c/b$ necessarily permits sufficiently weak under-cooperation.

To establish sufficiency for the entire family, for $0<u\leq1$ compare the mapping with its counterpart without recovery states:
\begin{equation}
\frac{\mathcal R_{k,\ell,\alpha}(u)}{\mathcal R_{k,0,\alpha}(u)}
=\frac{u^\ell}{u^\ell+(1-u)^k(1-u^\ell)}\leq1.
\label{eq:si-readout-recovery-bound}
\end{equation}
Both mappings equal zero at $u=0$. Concavity of $\mathcal R_{k,0,\alpha}$ and its tangent at $u=1$ therefore give
\begin{equation}
\mathcal R_{k,\ell,\alpha}(u)
\leq\mathcal R_{k,0,\alpha}(u)
\leq1-(1-\alpha)(1-u).
\label{eq:si-readout-linear-bound}
\end{equation}
The mutant has $u_{\mathrm X}=1-\delta$, so
\begin{equation}
\Lambda_{\mathrm{iid}}(\delta)
=\delta-\frac bc[1-\mathcal R_{k,\ell,\alpha}(1-\delta)]
\leq\delta\left[1-\frac bc(1-\alpha)\right].
\label{eq:si-readout-errorfree-payoff-bound}
\end{equation}
Combining necessity and sufficiency proves
\begin{equation}
\boxed{\Lambda_{\mathrm{iid}}(\delta)\leq0\ \text{for all }0<\delta\leq1
\quad\Longleftrightarrow\quad\alpha\leq1-\frac cb.}
\label{eq:si-readout-errorfree-threshold}
\end{equation}
This is a conditional error-free result for the independent under-cooperation family. Equality may include neutral mutants. The proof does not establish attraction from arbitrary reputation initializations, general evolutionary stability or resistance to temporally correlated mutants. In particular, $\alpha=0.8$ at $b/c=5$ is an analytical benchmark, rather than an imposed exact threshold for the finite-error figures.

\subsubsection{Finite-error security and the stationary branch limit}
\label{sec:si-readout-branches}

With finite errors, the resident assessment probability is below one, and both the assessment contrast and mapping slope change with readout. For configurations satisfying Eq.~\eqref{eq:si-readout-local-global-equivalence}, an interior security boundary is characterized by
\begin{equation}
\begin{aligned}
r&=\mathcal R_{\alpha,\beta}[\Phi_{\mathrm R}(r)],\\
1&=\frac bc\Delta\Phi(r)\mathcal R'_{\alpha,\beta}[\Phi_{\mathrm R}(r)].
\end{aligned}
\label{eq:si-readout-security-equations}
\end{equation}
These equations give an analytical characterization, generally requiring numerical root finding. For other configurations, Eq.~\eqref{eq:si-readout-global-condition} is the defining condition; an interior minimum can also determine the boundary. We did not impose a single connected security interval before evaluating the criterion.

For fixed norm, depths, $\beta$ and errors, define
\begin{equation}
F(r,\alpha;\beta)=\mathcal R_{\alpha,\beta}[\Phi_{\mathrm R}(r)]-r.
\label{eq:si-readout-branch-residual}
\end{equation}
The lower existence boundary $\alpha_{\mathrm{rep}}$ of the high-reputation branch is a non-degenerate stationary fold satisfying
\begin{equation}
\begin{aligned}
F(r,\alpha;\beta)&=0,\qquad \partial_rF(r,\alpha;\beta)=0,\\
\partial_\alpha F(r,\alpha;\beta)&\ne0,\qquad
\partial_{rr}F(r,\alpha;\beta)\ne0.
\end{aligned}
\label{eq:si-readout-fold-conditions}
\end{equation}
The first two conditions are equivalently $r=\mathcal R_{\alpha,\beta}[\Phi_{\mathrm R}(r)]$ and $\mathcal R'_{\alpha,\beta}[\Phi_{\mathrm R}(r)]\Phi'_{\mathrm R}(r)=1$. No $b/c$ occurs in these equations. Hence changing the payoff ratio can move $\alpha_{\mathrm{sec}}$ without moving $\alpha_{\mathrm{rep}}$. A stationary fold is not, by itself, a proof of the complete latent-state dynamics' stability; the distinction and full-state checks are described in Supplementary Note~\ref{sec:si-full-state-dynamics}.

Table~\ref{tab:si-readout-boundaries} gives the values underlying Extended Data Fig.~\ref{fig:ed-alpha-invasion}. At $b/c=5$ and $\beta=1$, the high-branch secure intervals are
\begin{equation}
\begin{aligned}
\text{Shunning:}\quad&0.59116052837<\alpha\leq0.80646444140,\\
\text{Stern Judging:}\quad&0.56698584543<\alpha\leq0.80360397146.
\end{aligned}
\label{eq:si-readout-feasible-intervals}
\end{equation}
For these two representative configurations, below $\alpha_{\mathrm{rep}}$ the high branch is absent; at equality the two stationary roots coalesce at the fold. Above that endpoint and up to $\alpha_{\mathrm{sec}}$, the selected high branch excludes the tested family; above $\alpha_{\mathrm{sec}}$, it permits under-cooperation invasion. This classification is specific to the displayed interval structure. Neither Scoring nor Simple Standing has the corresponding high-branch fold in this scan. Their cooperation declines continuously when $\alpha$ is lowered, reaching approximately $0.338$ and $0.622$ at $\alpha=0$. Thus the absence of a positive invasion exponent cannot be interpreted as preservation of high cooperation.

\begin{table}[htbp]
\centering
\caption{\textbf{Readout boundaries for the representative configurations.}
Parameters are $e_1=e_2=0.02$ and $\beta=1$; security boundaries use $b/c=5$. ``None'' denotes absence of the corresponding stationary high-branch fold, not a fold at $\alpha=0$. The value $\alpha_{\mathrm{sec}}=1$ for Simple Standing denotes security at the allowed endpoint. The cooperation column is the intended cooperation of the selected pure-DISC resident at $\alpha=0.8$; the final column is the critical payoff ratio for excluding the complete independent under-cooperation family at that readout.}
\label{tab:si-readout-boundaries}
\begin{tabular}{lccccc}
\toprule
Norm & $(k_1,k_2)$ & $\alpha_{\mathrm{rep}}$ & $\alpha_{\mathrm{sec}}$ & $C_{\mathrm{DISC}}(0.8)$ & $(b/c)_{\mathrm{crit}}(0.8)$\\
\midrule
Scoring & $(2,0)$ & None & $0.807064$ & $0.989782$ & $4.827195$\\
Shunning & $(3,0)$ & $0.591161$ & $0.806464$ & $0.986341$ & $4.835173$\\
Stern Judging & $(4,3)$ & $0.566986$ & $0.803604$ & $0.986509$ & $4.910131$\\
Simple Standing & $(1,2)$ & None & $1$ & $0.977886$ & $1.955313$\\
\bottomrule
\end{tabular}
\end{table}

The fold and security conditions also explain where the two boundaries meet. Under Shunning, let $a=(1-e_1)(1-2e_2)$. The effective assessment probabilities give $\Phi_{\mathrm R}(r)=e_2+ar^2$ and $\Delta\Phi(r)=ar$. At the fold, $\mathcal R'\Phi'_{\mathrm R}=1$, hence $\Delta\Phi\mathcal R'=1/2$. Because $k_2=0$ makes local and complete-family criteria equivalent, the critical payoff ratio at this fold is exactly $b/c=2$. Under Stern Judging,
\begin{equation}
\Delta\Phi(r)=a(2r-1),\qquad
\Phi'_{\mathrm R}(r)=2\Delta\Phi(r)-e_1(1-2e_2).
\end{equation}
The local critical ratio at its fold is therefore
\begin{equation}
\left(\frac bc\right)_{\mathrm{fold,loc}}
=\frac{\Phi'_{\mathrm R}(r_{\mathrm{fold}})}{\Delta\Phi(r_{\mathrm{fold}})}
=2-\frac{e_1(1-2e_2)}{\Delta\Phi(r_{\mathrm{fold}})}
\simeq1.9709729883.
\label{eq:si-readout-fold-critical-bc}
\end{equation}
The complete-family calculation agreed with this local value at the displayed Stern Judging fold. These results account for the horizontal branch limits and the payoff-dependent security boundaries in Extended Data Fig.~\ref{fig:ed-alpha-invasion}j,k.

\subsubsection{Distinct effects of warning and recovery readouts}
\label{sec:si-readout-beta}

Readout parameters change how evidence affects current treatment without changing the latent transition rule at a fixed assessment probability. Let
$Z(u)=(1-u)^{k_1}+u^{k_2}-(1-u)^{k_1}u^{k_2}$.
Differentiating Eq.~\eqref{eq:si-reputation-transform} at fixed $u$ gives
\begin{equation}
\begin{aligned}
\partial_\alpha\mathcal R_{\alpha,\beta}(u)
&=\frac{u^{k_2+1}}{Z(u)}
\sum_{m=1}^{k_1}m\alpha^{m-1}(1-u)^m\geq0,\\
\partial_\beta\mathcal R_{\alpha,\beta}(u)
&=-\frac{(1-u)^{k_1+1}}{Z(u)}
\sum_{n=1}^{k_2}n\beta^{n-1}u^n\leq0.
\end{aligned}
\label{eq:si-readout-parameter-derivatives}
\end{equation}
The sums are empty when the corresponding path is absent. Thus $\beta$ has exactly no effect for $k_2=0$, as in the representative Scoring and Shunning configurations. At a regular self-consistent branch, implicit differentiation gives, for $\theta\in\{\alpha,\beta\}$,
\begin{equation}
\frac{\partial r^*}{\partial\theta}
=\frac{\partial_\theta\mathcal R_{\alpha,\beta}(u^*)}
{1-\mathcal R'_{\alpha,\beta}(u^*)\Phi'_{\mathrm R}(r^*)}.
\label{eq:si-readout-branch-sensitivity}
\end{equation}
On regular segments with positive denominator, increasing $\alpha$ increases resident reputation and increasing $\beta$ decreases it. The denominator vanishes at a fold, where ordinary branch sensitivity fails. Changes in invasion fitness additionally involve the mutant mapping and the changed resident assessment contrast; the signs in Eq.~\eqref{eq:si-readout-parameter-derivatives} alone do not prove monotonic evolutionary security.

Near $u=1$, for $k_1\geq1$ and arbitrary $\beta$, the leading mapping response remains $\mathcal R_{\alpha,\beta}(1-v)=1-(1-\alpha)v+O(v^2)$. The direct $\beta$ dependence is of order $v^{k_1+1}$ in reputation, because recovery states require prior loss of established good standing. At full warning readout $\alpha=1$, this dependence is made explicit by
\begin{equation}
\begin{aligned}
1-\mathcal R_{1,\beta}(1-v)
&=\left(1+\sum_{n=1}^{k_2}\beta^n\right)v^{k_1+1}
+O(v^{k_1+2}),\\
\mathcal R'_{1,\beta}(1-v)
&=(k_1+1)\left(1+\sum_{n=1}^{k_2}\beta^n\right)v^{k_1}
+O(v^{k_1+1}).
\end{aligned}
\label{eq:si-readout-beta-asymptotic}
\end{equation}
When $\beta=1$, the coefficient is $k_2+1$. Deep warning paths therefore suppress the marginal response even when recovery strictness is maximal. In contrast, lowering $\alpha$ introduces a first-order response through the first warning. The distinction concerns the high-reputation expansion; at finite deviations and finite errors, both parameters also alter the resident environment.

These expressions explain the different representative outcomes. At $\alpha=\beta=1$, the Scoring, Shunning and Stern Judging responses scale as $v^2$, $v^3$ and $v^4$, whereas Simple Standing $(1,2)$ retains a response of approximately $6v$. At the finite errors used here, Simple Standing has $G_{\mathrm{global}}\simeq1.342945$ even at $\alpha=1$. This does not contradict the conditional error-free threshold, whose resident has $u=1$ exactly.

We also evaluated recovery readout at $\alpha=1$, $b/c=5$ and $e_1=e_2=0.02$. In Stern Judging $(4,3)$, increasing $\beta$ from zero to one reduced the invasion-window width from $0.431865$ to $0.365299$, but the local derivative remained positive throughout this range (approximately $0.999760$--$0.999936$). Thus recovery strictness narrowed finite exploitation without removing near-resident invasion in this configuration. In Simple Standing $(1,2)$, the window at $\beta=0$ had width $0.090320$ and disappeared at
\begin{equation}
\beta_{\mathrm{sec}}\simeq0.7408874073,\qquad\alpha=1.
\label{eq:si-readout-beta-sec}
\end{equation}
At and above this boundary the complete tested family had non-positive invasion fitness. These numerical statements are conditional on the stated norm, depths and errors. In particular, the ability to restore security using $\beta$ alone is not a generic property of Simple Standing at every depth: with the one-step structure $(1,1)$, none of the four norms excluded the complete family at $\alpha=1$ even when $\beta=1$.

\subsection{Temporally structured opportunism}
\label{sec:si-temporal-opportunism}

\subsubsection{Markov-switching opportunism}

Let $Z_a\in\{0,1\}$ be the mode at donor action $a$: $0$ follows DISC and $1$ follows ALLD. At the end of the action the mode changes with transition matrix
\begin{equation}
T_{\mathrm a}=\begin{pmatrix}
1-\delta(1-\eta_{\mathrm a}) & \delta(1-\eta_{\mathrm a})\\
(1-\delta)(1-\eta_{\mathrm a}) & 1-(1-\delta)(1-\eta_{\mathrm a})
\end{pmatrix}
=\eta_{\mathrm a}I+(1-\eta_{\mathrm a})\mathbf 1(1-\delta,\delta),
\label{eq:si-mode-transition-action}
\end{equation}
where $0<\delta<1$ and $0\leq\eta_{\mathrm a}<1$. Direct multiplication gives the stationary mode distribution $(1-\delta,\delta)$. The non-unit eigenvalue is $\eta_{\mathrm a}$, hence $\operatorname{Corr}(Z_a,Z_{a+n})=\eta_{\mathrm a}^{n}$ at stationarity. The lengths of DISC and ALLD runs are geometric, with means $1/[\delta(1-\eta_{\mathrm a})]$ and $1/[(1-\delta)(1-\eta_{\mathrm a})]$, respectively. Thus $\delta$ controls the long-run amount of exploitation and $\eta_{\mathrm a}$ controls its temporal arrangement.

A given observer need not see every donor action. If each action updates that observer's focal record independently with probability $p_{\mathrm{obs}}$, the number $G$ of action intervals between consecutive observations is geometric: $\Pr(G=g)=p_{\mathrm{obs}}(1-p_{\mathrm{obs}})^{g-1}$ for $g\geq1$. Averaging $T_{\mathrm a}^{G}$ gives the mode kernel on the observation clock,
\begin{equation}
T_{\mathrm{obs}}=\eta I+(1-\eta)\mathbf 1(1-\delta,\delta),
\qquad
\eta=\mathbb E[\eta_{\mathrm a}^{G}]
=\frac{p_{\mathrm{obs}}\eta_{\mathrm a}}{1-(1-p_{\mathrm{obs}})\eta_{\mathrm a}}.
\label{eq:si-observed-persistence}
\end{equation}
All displayed temporal results use this observed persistence $\eta$. For the finite-population observation convention used to calibrate the clock, one rare mutant is observed by any specified resident with probability $p_{\mathrm{obs}}=1/(N-1)=1/149$ per mutant donor action. A different observation probability changes the action-time value $\eta_{\mathrm a}$ corresponding to the same $\eta$, but not the joint chain evaluated on the observation clock. The observed mean DISC and ALLD run lengths are $1/[\delta(1-\eta)]$ and $1/[(1-\delta)(1-\eta)]$ observation intervals.

Let $M_0=\mathbf M(\Phi_{\mathrm R}^*)$ and $M_1=\mathbf M(\Phi_{\mathrm D}^*)$ be the latent-reputation transition matrices after an observed DISC-mode and ALLD-mode action. The assessment is drawn and the focal reputation updated under the current mode, then the mode switches before the next observed action. Ordering joint states as $(0,L)$ followed by $(1,L)$ gives
\begin{equation}
J(\delta,\eta)=
\begin{pmatrix}
[1-\delta(1-\eta)]M_0 & \delta(1-\eta)M_0\\
(1-\delta)(1-\eta)M_1 & [1-(1-\delta)(1-\eta)]M_1
\end{pmatrix}.
\label{eq:si-temporal-joint-kernel}
\end{equation}
If $vJ=v$, $v\mathbf1=1$, and $\boldsymbol\omega$ is the latent-state good-readout vector of Eq.~\eqref{eq:si-full-state-readout}, then
\begin{equation}
r_{\mathrm X}^*(\delta,\eta)
=v\begin{pmatrix}\boldsymbol\omega\\\boldsymbol\omega\end{pmatrix},
\qquad
\Lambda(\delta,\eta)=\delta r_{\mathrm R}^*-\frac bc[r_{\mathrm R}^*-r_{\mathrm X}^*(\delta,\eta)].
\label{eq:si-temporal-invasion}
\end{equation}
The stationary mode mass in $v$ is $(1-\delta,\delta)$, but mode and latent reputation generally remain correlated. For example, after several exploitative observations the record is more likely to be in a warning or bad state than after several DISC observations. Consequently, replacing the process by the scalar map $\mathcal R[(1-\delta)\Phi_{\mathrm R}^*+\delta\Phi_{\mathrm D}^*]$ is generally incorrect when $\eta>0$.

At $\eta=0$, each new mode is independent of the previous mode and latent state. Marginalizing $J$ then yields the reputation kernel $(1-\delta)M_0+\delta M_1=\mathbf M(\Phi_{\mathrm R}^*-\delta r_{\mathrm R}^*\Delta\Phi^*)$, because $\mathbf M$ is affine in the positive-assessment probability. Thus Eq.~\eqref{eq:si-temporal-invasion} reduces exactly to Eq.~\eqref{eq:si-iid-invasion}. At $\delta=0$ or $1$, one mode is absorbing and the stationary reputations are $r_{\mathrm R}^*$ and $r_{\mathrm D}^*$, independently of $\eta<1$.

For fixed $0<\delta<1$, as $\eta\to1^{-}$ the mean duration of each mode diverges. If both mode-specific latent chains are ergodic, their records relax during each long episode, so
\begin{equation}
\lim_{\eta\to1^-}r_{\mathrm X}^*(\delta,\eta)
=(1-\delta)r_{\mathrm R}^*+\delta r_{\mathrm D}^*,
\qquad
\lim_{\eta\to1^-}\Lambda(\delta,\eta)
=\delta\left\{r_{\mathrm R}^*-\frac bc(r_{\mathrm R}^*-r_{\mathrm D}^*)\right\}
=\delta\Lambda_{\mathrm{ALLD}}.
\label{eq:si-temporal-persistent-limit}
\end{equation}
At $\eta=1$ itself, the behavioural chain has two closed mode classes and has no unique stationary mode mixture; the preceding expression is a one-sided limit, not an ergodic stationary formula at that endpoint. When the iid exponent is positive at a fixed $\delta$ but $\Lambda_{\mathrm{ALLD}}<0$, continuity ensures at least one boundary at an intermediate $\eta$. It does not imply pointwise monotonicity or equivalence between a local $\delta\downarrow0$ boundary and closure of every finite-$\delta$ invasion interval. Those boundaries were checked separately in the numerical scans.

\subsubsection{Periodic exploitation and reputation washing}

A deterministic counterpart alternates between $h$ consecutive DISC donor actions and $\ell$ consecutive ALLD donor actions.
The exploitative fraction is
\begin{equation}
\delta
=
\frac{\ell}{h+\ell}.
\end{equation}

Let $\boldsymbol{\pi}_0$ be the internal-state distribution at the beginning of the DISC block.
The periodic stationary condition is
\begin{equation}
\boldsymbol{\pi}_0
=
\boldsymbol{\pi}_0
M_0^hM_1^\ell,
\qquad
\boldsymbol{\pi}_0\mathbf{1}=1.
\end{equation}
If $\boldsymbol{\pi}_t$ denotes the internal-state distribution immediately before donor action $t$ within the cycle, the cycle-averaged good reputation is
\begin{equation}
\bar r_P(h,\ell)
=
\frac{1}{h+\ell}
\sum_{t=0}^{h+\ell-1}
\boldsymbol{\pi}_t\boldsymbol{\omega}.
\label{eq:SI_periodic_reputation}
\end{equation}

The mutant intends cooperation only during the DISC block, so
\begin{equation}
q_\mathrm P
=
\frac{h}{h+\ell}r_\mathrm R^*,
\end{equation}
and again saves
\begin{equation}
S_\mathrm P
=
\delta r_\mathrm R^*
\end{equation}
relative to resident DISC.
Its invasion exponent is therefore
\begin{equation}
\frac{\lambda_\mathrm P(h,\ell)}{(1-e_1)c}
=
\delta r_\mathrm R^*
-
\frac{b}{c}
\left[
r_\mathrm R^*-\bar r_\mathrm P(h,\ell)
\right].
\label{eq:SI_periodic_invasion}
\end{equation}
Whenever $\bar r_\mathrm P<r_\mathrm R^*$, the corresponding lower threshold is
\begin{equation}
R_\mathrm P(h,\ell)
=
\frac{\delta r_\mathrm R^*}
{r_\mathrm R^*-\bar r_\mathrm P(h,\ell)}.
\label{eq:SI_periodic_threshold}
\end{equation}

If the block lengths are increased while their ratio is kept fixed, each behavioural episode becomes long enough for reputation to approach its mode-specific equilibrium.
In this limit,
\begin{equation}
R_\mathrm P(nh,n\ell)
\rightarrow
R_\mathrm L
\qquad
(n\rightarrow\infty).
\end{equation}
Any additional advantage of a finite washing cycle therefore arises from repeatedly exploiting transient intermediate reputation states.

\subsubsection{Reputation recovery after finite exploitation}

Periodic exploitation should be distinguished from genuine behavioural reform.
Let $\boldsymbol{\pi}_R$ denote the stationary internal-state distribution under resident DISC behaviour,
\begin{equation}
\boldsymbol{\pi}_\mathrm RM_0
=
\boldsymbol{\pi}_\mathrm R,
\qquad
\boldsymbol{\pi}_\mathrm R\boldsymbol{\omega}
=
r_\mathrm R^*.
\end{equation}
After $\ell$ consecutive exploitative donor actions followed by $h$ consecutive DISC donor actions, the good-reputation probability is
\begin{equation}
r_{\mathrm{rec}}(\ell,h)
=
\boldsymbol{\pi}_\mathrm R
M_1^\ell
M_0^h
\boldsymbol{\omega}.
\label{eq:SI_recovery}
\end{equation}
If $M_0$ has a unique stationary distribution,
\begin{equation}
\lim_{h\to\infty}
r_{\mathrm{rec}}(\ell,h)
=
r_\mathrm R^*.
\end{equation}
A finite exploitative history is therefore not permanently irreversible.
Repeated reputation washing differs from genuine reform because every new exploitative episode again moves the reputation away from the resident distribution.

\subsection{Interpretation in terms of accumulated evidence}

The unified invasion condition,
\begin{equation}
\frac{\lambda_\mathrm X}{(1-e_1)c}
=
S_\mathrm X-\frac{b}{c}D_\mathrm X,
\end{equation}
shows that resistance to invasion depends on how strongly a behavioural deviation is translated into reputational loss.

Three components contribute to this translation.

First, the social norm determines the assessment contrast $\Delta\Phi^*$, and therefore how differently cooperation and defection are evaluated in the resident environment.

Second, the inertia depths $k_1$ and $k_2$ determine how much consistent evidence is required to overturn an established reputation.
Under the direct-reset rule, moving from $G$ to $B$ requires $k_1+1$ consecutive negative assessments, whereas moving from $B$ to $G$ requires $k_2+1$ consecutive positive assessments.

Third, the readout parameters $\alpha$ and $\beta$ determine how intermediate evidence affects current behaviour before an endpoint reputation is reached.
Thus, evidence accumulation and current trust are controlled separately.

Under full readout inertia, $\alpha=\beta=1$, the waiting time for $k_1+1$ consecutive negative assessments with per-assessment probability $q$ is
\begin{equation}
\mathbb{E}[\tau_{\mathrm G\rightarrow \mathrm B}]
=
\frac{1-q^{k_1+1}}
{(1-q)q^{k_1+1}},
\end{equation}
whereas the waiting time for $k_2+1$ consecutive positive assessments with probability $p$ is
\begin{equation}
\mathbb{E}[\tau_{\mathrm B\rightarrow \mathrm G}]
=
\frac{1-p^{k_2+1}}
{(1-p)p^{k_2+1}}.
\end{equation}

These expressions provide a simple interpretation of the inertia depths, while the invasion calculations above use the complete reputation processes.
Together, the results distinguish the amount of evidence required to revise a reputation from the way partially accumulated evidence is used in current interactions.

\subsection{Observation aggregation and quantitative reputation benchmarks}
\label{sec:si-benchmark-mechanisms}

We compare alternative reputation-processing rules in the same donation game, four second-order norms, implementation and assessment errors, pure-DISC resident environment and rare-mutant payoff identity as above. Only the conversion of assessments to a good readout changes. The original studies employ different strategy sets and observation procedures, so these calculations are controlled adaptations, not replications of their complete evolutionary outcomes~\cite{si-michelmata2024,si-schmid2023}. We select the resident root reached from the all-good stationary closure and report results conditional on that branch. In particular, a fixed-point root alone does not establish that every microscopic initialization selects it.

\subsubsection{Aggregation of sampled interactions}
\label{sec:si-aggregation}

Michel-Mata \textit{et al.} define an aggregating discriminator by the number $M\geq1$ of sampled donor interactions and a tolerance threshold $q_{\mathrm A}\in[0,1]$ (denoted $q$ in the source paper)~\cite{si-michelmata2024}. An observer samples interactions uniformly with replacement from the donor's available interaction pool, assesses each under its norm and assigns a binary good reputation when at least $Q=\lceil q_{\mathrm A}M\rceil$ of the $M$ assessments are positive. We write $Q$ explicitly because the threshold is a count. When $q_{\mathrm A}=0$, $Q=0$ makes every readout good; for $0<q_{\mathrm A}\leq1$, $1\leq Q\leq M$. The representative published ``look twice, forgive once'' rule has $(M,q_{\mathrm A},Q)=(2,1/2,1)$.

Suppose each sampled interaction has an independent positive assessment with probability $u$. The number $K$ of positive assessments is binomial, since each particular sequence containing $j$ positives has probability $u^j(1-u)^{M-j}$ and there are ${M\choose j}$ such sequences. Therefore the exact stationary-pool good-readout map is
\begin{equation}
F_{\mathrm A}(u;M,Q)
=\Pr(K\geq Q)
=\sum_{j=Q}^{M}{M\choose j}u^j(1-u)^{M-j}.
\label{eq:si-aggregation-map}
\end{equation}
In particular, $F_{\mathrm A}(u;2,1)=1-(1-u)^2=2u-u^2$, $F_{\mathrm A}(u;M,M)=u^M$ and $F_{\mathrm A}(u;M,1)=1-(1-u)^M$. The expression gives $F_{\mathrm A}(0)=0$ and $F_{\mathrm A}(1)=1$ for $1\leq Q\leq M$. Its derivative on $0<u<1$ is $M{M-1\choose Q-1}u^{Q-1}(1-u)^{M-Q}$, obtained by differentiating the sum and cancelling adjacent terms.

Within the present stationary private-assessment closure, an all-DISC resident satisfies
\begin{equation}
r_{\mathrm R}^*=F_{\mathrm A}[\Phi_{\mathrm R}(r_{\mathrm R}^*);M,Q],
\qquad
r_{\mathrm X}^*(\delta)=F_{\mathrm A}(\Phi_{\mathrm R}^*-\delta r_{\mathrm R}^*\Delta\Phi^*;M,Q).
\label{eq:si-aggregation-resident-mutant}
\end{equation}
The rare-mutant exponent is Eq.~\eqref{eq:si-iid-invasion} with $\mathcal R$ replaced by $F_{\mathrm A}$. For $M=2,Q=1$, write $u_{\mathrm R}=\Phi_{\mathrm R}^*$ and $\Delta=\Delta\Phi^*$. Expanding the quadratic map and using $r_{\mathrm R}^*=F_{\mathrm A}(u_{\mathrm R})$ gives
\begin{equation}
\Lambda_{\mathrm A}(\delta)
=r_{\mathrm R}^*\delta\left[1-2\frac bc(1-u_{\mathrm R})\Delta
-\frac bc r_{\mathrm R}^*\Delta^2\delta\right].
\label{eq:si-aggregation-quadratic}
\end{equation}
For the Stern Judging parameters of Extended Data Fig.~\ref{fig:ed-aggregation-quantitative}, this is $\Lambda_{\mathrm A}(\delta)=0.59987648\delta-4.37759510\delta^2$, with the positive interval ending at $\delta=0.13703334$. In the error-free all-good limit it becomes $\delta-(b/c)\delta^2$, so the boundary is $\delta=c/b$. The original paper's probabilistic cooperator uses an unconditional cooperation probability $p$, rather than our conditional omission probability $\delta$. Their coordinates are equivalent as $p=1-\delta$ in this error-free all-good limit. Under Scoring the assessment also ignores the recipient's reputation, permitting an equivalent mean-field parameterization $p=r_{\mathrm R}^*(1-\delta)$; this equivalence does not extend to the other norms with private disagreement and errors.

Equation~\eqref{eq:si-aggregation-map} requires independent \emph{sampled} assessments. The benchmark follows the original random-with-replacement protocol in the stationary infinite-interaction-pool limit. Each draw then sees the marginal positive-assessment probability $u$, even if actions were temporally correlated before entering the pool; independently redrawing $M$ times makes Eq.~\eqref{eq:si-aggregation-map} valid and the benchmark independent of $\eta$. In a finite pool with $L$ recorded interactions and $K$ positive entries, the conditional probability is $F_{\mathrm A}(K/L;M,Q)$, but averaging over a random, temporally correlated $K$ gives $\mathbb E[F_{\mathrm A}(K/L)]$, generally not $F_{\mathrm A}(\mathbb E[K/L])$. A rule based on the last $M$ consecutive assessments is different again. Its exact state is the pair of behavioural mode $z$ and assessment word $w\in\{0,1\}^{M}$, with transitions
\begin{equation}
\Pr[(z,w)\to(z',\operatorname{shift}(w,y))]
=T_{\mathrm{obs}}(z,z')\,u_z^{y}(1-u_z)^{1-y},
\quad y\in\{0,1\},
\label{eq:si-aggregation-window-kernel}
\end{equation}
where $u_0=\Phi_{\mathrm R}^*$ and $u_1=\Phi_{\mathrm D}^*$. Its good readout is $\mathbf1\{\sum w_i\geq Q\}$ and must be averaged against the joint stationary distribution. We do not use this consecutive-window rule in the Extended Data benchmark. The stationary infinite-pool limit and $\eta\to1^{-}$ limit need not commute for finite records; consequently the Aggregation landscape in Extended Data Fig.~\ref{fig:ed-aggregation-quantitative}c need not approach $\delta\Lambda_{\mathrm{ALLD}}$.

The displayed Aggregation family uses both two-observation thresholds, all four count thresholds for $M=4$, and the source paper's $q_{\mathrm A}=1/2$ depth series at $M=6,8,10$:
\[
(M,Q)\in\{(2,1),(2,2),(4,1),(4,2),(4,3),(4,4),
(6,3),(8,4),(10,5)\}.
\]
The coloured curve is $(2,1)$; the others are grey. These are parameter examples allowed or analysed in the original paper, not the full space of aggregation rules. Some strict cases select a lower-reputation resident branch or fail even the ALLD test under the present adaptation, so the representative result must not be generalized to the whole family.

\subsubsection{Accumulated quantitative scores}
\label{sec:si-quantitative}

Schmid \textit{et al.} retain an observer-specific integer reputation score $z\in\{-R,-R+1,\ldots,R\}$~\cite{si-schmid2023}. A positive assessment increments the score to $\min(z+1,R)$, and a negative assessment decrements it to $\max(z-1,-R)$. The score is read as good precisely when $z\geq S$. The symmetric score range with $S=0$ is the paper's principal frame; $R=5,S=0$ is the representative Extended Data setting, and the grey curves cover its publicly examined frames $R=1,2,3,4$ at $S=0$.

For a constant positive-assessment probability $0<u<1$, the score is a reflecting birth--death chain. Index states by $j=z+R\in\{0,\ldots,2R\}$. Each interior upward move has probability $u$, each downward move probability $1-u$, and attempted moves beyond the endpoints remain at the endpoint. Stationary detailed balance between adjacent states requires $\pi_j u=\pi_{j+1}(1-u)$. With $t=u/(1-u)$ and normalization, the unique stationary distribution is
\begin{equation}
\pi_j(u)=\frac{t^j}{\sum_{\ell=0}^{2R}t^\ell},
\qquad j=0,\ldots,2R.
\label{eq:si-quantitative-stationary}
\end{equation}
Summing the stationary mass of states $z\geq S$ proves the quantitative good-readout map
\begin{equation}
F_{\mathrm Q}(u;R,S)
=\frac{\sum_{j=S+R}^{2R}t^j}{\sum_{j=0}^{2R}t^j},
\qquad -R\leq S\leq R.
\label{eq:si-quantitative-map}
\end{equation}
At $u=1/2$, the distribution is uniform and $F_{\mathrm Q}=(R-S+1)/(2R+1)$; at $u=0$ or $1$ it concentrates at $-R$ or $R$, respectively. For $R=1,S=0$, Eq.~\eqref{eq:si-quantitative-map} reduces to $u/(1-u+u^2)$, providing a short independent check. The stationary resident and iid mutant equations are Eq.~\eqref{eq:si-aggregation-resident-mutant} with $F_{\mathrm A}$ replaced by $F_{\mathrm Q}$; the payoff remains Eq.~\eqref{eq:si-iid-invasion} with the same replacement.

A temporally persistent quantitative mutant cannot be computed by substituting its mean assessment probability into $F_{\mathrm Q}$. Let $K_0=K(\Phi_{\mathrm R}^*)$ and $K_1=K(\Phi_{\mathrm D}^*)$ be the reflecting score kernels and let $h_j=\mathbf1\{j-R\geq S\}$. The exact observation-clock joint transition is
\begin{equation}
J_{\mathrm Q}=\begin{pmatrix}
T_{\mathrm{obs}}(0,0)K_0&T_{\mathrm{obs}}(0,1)K_0\\
T_{\mathrm{obs}}(1,0)K_1&T_{\mathrm{obs}}(1,1)K_1
\end{pmatrix},
\qquad
r_{\mathrm X}^*=v_{\mathrm Q}\begin{pmatrix}h\\h\end{pmatrix},
\label{eq:si-quantitative-joint}
\end{equation}
where $v_{\mathrm Q}J_{\mathrm Q}=v_{\mathrm Q}$ and $v_{\mathrm Q}\mathbf1=1$. This has $2(2R+1)$ states. At $\eta=0$ it reduces to the iid map. At $\delta=0$ and $1$ it gives the resident and ALLD score equilibria. For $0<\delta<1$, its $\eta\to1^{-}$ limit is the weighted mixture of those equilibria, and Eq.~\eqref{eq:si-temporal-persistent-limit} holds. These limits and stationary residuals were checked against the numerical output.

The quantitative source model also varies observation probability and uses leading-eight norms, its own error convention, initial score $z=0$, and finite-population selection--mutation dynamics. Here we transplant only its score-update and threshold-readout rule into the four second-order norms and payoff framework of Fig.~\ref{fig:strategy-invasion}. The benchmark selects the high resident root by starting the stationary closure at $r=1$; this choice differs from the source paper's microscopic $z=0$ initialization, and Scoring and Shunning have additional lower roots for $R=5$. Results are therefore conditional on the reported branch. In that branch, $1-r_{\mathrm R}^*\simeq4.61\times10^{-9}$ at $R=5,S=0$. The four norms differ only in assessments involving recipients read as bad, so their representative invasion curves nearly coincide, with a maximum difference $6.19\times10^{-8}$ in $\Lambda$ over the calculated grid. The near-coincidence is branch- and parameter-specific; the $R=1$ family retains visible norm differences. Under Stern Judging, the representative iid interval ends at $\delta=0.383206$, and the observed-persistence boundary is approximately $\eta=0.664552$. These boundaries characterize the tested mutant family, not general evolutionary stability.

\subsection{Numerical strategy dynamics and parameter scans}

For every strategy composition, we first integrated the reputation subsystem from $r_I(0)=1$ to a stable equilibrium and then evaluated the payoffs. We integrated the resulting replicator dynamics from $M=512$ common initial compositions drawn uniformly from the simplex with a scrambled Sobol sequence. If $\boldsymbol f_\infty^{(m)}$ is the terminal strategy composition reached from initial point $m$, basin-weighted intended cooperation is
\begin{equation}
\left\langle\rho_{\mathrm{int}}\right\rangle_{\Delta}
=\frac{1}{M}\sum_{m=1}^{M}\rho_{\mathrm{int}}\!\left(\boldsymbol f_\infty^{(m)}\right).
\label{eq:si-basin-weighted-cooperation}
\end{equation}
The same initial compositions were used for all norms and parameter values. The shaded regions in the simplex panels use a separate binary classification. An initial composition belongs to the cooperative basin if its terminal intended cooperation is at least $0.5$.

Basin-weighted cooperation combines two contributions. The first is the fraction of the simplex occupied by basins leading to cooperative states. The second is the cooperation attained inside those basins. It can therefore be high because a large set of initial conditions reaches cooperation, because the cooperative attractor itself has a high cooperation level, or because of both effects. The local thresholds answer a different question and do not determine basin size.

For the one-step scan, we fixed $k_1=k_2=1$ and evaluated the $(\alpha,\beta)$ plane at $b/c=5$. We additionally computed simplex flows at $(\alpha,\beta)=(1,0)$ and $(0,1)$ to separate the effects of tolerance and strictness. At every grid point, Eq.~\eqref{eq:SI_general_invasion} was evaluated for $\mathrm{ALLD}$ and $\mathrm{ALLC}$ to classify whether $\mathrm{ALLD}$ invades, whether $\mathrm{ALLD}$ is excluded but $\mathrm{ALLC}$ invades, or whether both mutants are excluded.

For the benefit-to-cost scan, we integrated the strategy dynamics over the same initial compositions at each $b/c$. Curve segments were classified independently by the two pure-$\mathrm{DISC}$ invasion exponents. This permits similar cooperation levels arising from different equilibrium structures to be distinguished.

For the error scan, we fixed $b/c=5$ and evaluated the selected depth pair for each norm over the $(e_1,e_2)$ plane. At every point, the resident reputation was selected from the all-good initial condition before the two rare-mutant growth rates were evaluated. The resulting regions measure robustness of pure $\mathrm{DISC}$ to invasion and do not by themselves measure the cooperation level of a mixed attractor.

For the depth scan, we fixed $\alpha=\beta=1$ and evaluated all non-negative integer pairs $(k_1,k_2)$. We grouped allocations by total depth $K=k_1+k_2$ and defined the upper envelope as the largest observed basin-weighted cooperation at each $K$. The pairs $(2,0)$, $(3,0)$, $(4,3)$ and $(1,2)$ were selected beforehand from the finite-population recognition criteria in Supplementary Note~\ref{sec:si-finite-monte-carlo}; evolutionary payoffs were not used in their selection.



\section{Finite-population Monte Carlo simulations}
\label{sec:si-finite-monte-carlo}

\subsection{Purpose and simulated reputation records}

The population-level equations in Supplementary Note~\ref{sec:si-model} describe strategy-specific probabilities of a good reputation and close the private assessments of a donor and an observer by conditional independence. The finite-population simulations instead retain the internal reputation state of each simulated observer--target record. They therefore capture sampling fluctuations, correlations between private opinions, and differences among targets of the same strategy without imposing that closure.

We consider a fixed population of $N=n_{\mathrm{DISC}}+n_{\mathrm{ALLC}}+n_{\mathrm{ALLD}}$ individuals. Unless otherwise stated, $N=150$ and $n_{\mathrm{DISC}}=n_{\mathrm{ALLC}}=n_{\mathrm{ALLD}}=50$. Individuals retain their strategies throughout a run. The simulation stores
\[
\mathbf L=\{L_{vj}:v\in\mathcal V_{\mathrm{DISC}},\ 1\leq j\leq N\},
\]
where $\mathcal V_{\mathrm{DISC}}$ is the set of discriminators and $L_{vj}$ is the internal state held by discriminator $v$ for target $j$. The diagonal records $L_{vv}$ are included. Records held by $\mathrm{ALLC}$ and $\mathrm{ALLD}$ players are not simulated because those strategies do not consult reputations when selecting actions. This reduction does not alter any event in the fixed-strategy reputation process: every observer is sampled from $\mathcal V_{\mathrm{DISC}}$, and every reputation-dependent donor is also a member of that set. All retained records are initialized in state $\mathrm G$.

The internal state space, the readout function and the state-transition rules are exactly those in Supplementary Eqs.~\eqref{eq:si-binary-readout} and \eqref{eq:si-markov-matrix}. In particular, a consultation draws a new binary reputation $R_{vj}\sim\operatorname{Bernoulli}[\sigma(L_{vj})]$ without changing $L_{vj}$. Two consultations of the same record are conditionally independent given its current internal state. The internal state changes only when its observer assesses its target in the donor role.

\subsection{Monte Carlo event}
\label{sec:si-monte-carlo-event}

One event proceeds as follows.

\begin{enumerate}
\item Draw a donor $D$ and a recipient $R$ uniformly without replacement from the population. Draw an observer $V$ independently and uniformly from $\mathcal V_{\mathrm{DISC}}$. The observer may coincide with the donor or recipient.

\item Determine the donor's intended action. An $\mathrm{ALLC}$ donor intends to cooperate and an $\mathrm{ALLD}$ donor intends to defect. If $D$ uses $\mathrm{DISC}$, draw $R_{DR}\sim\operatorname{Bernoulli}[\sigma(L_{DR})]$ and let the donor intend to cooperate if and only if $R_{DR}=1$.

\item Apply implementation error. Intended cooperation becomes realized defection with probability $e_1$; intended defection remains defection.

\item Independently draw the observer's binary reputation of the recipient:
\[
R_{VR}\sim\operatorname{Bernoulli}[\sigma(L_{VR})].
\]
Apply the chosen social norm to the realized action and $R_{VR}$, and reverse the prescribed assessment with probability $e_2$.

\item Update the single internal record $L_{VD}$ using the transition for the recorded positive or negative assessment. All other internal records are left unchanged.
\end{enumerate}

Algorithm~\ref{alg:si-finite-reputation} gives the complete simulation logic. The notation $T^+(L)$ and $T^-(L)$ denotes the deterministic next internal state after a positive or negative assessment, respectively.

\begin{algorithm}[htbp]
\caption{Finite-population private-reputation simulation}
\label{alg:si-finite-reputation}
\begin{algorithmic}[1]
\Require Strategy sets, norm $(p,q)$, $e_1,e_2$, $\alpha,\beta$, $k_1,k_2$, number of sweeps $T$, burn-in $T_0$, sampling interval $\Delta T$
\State Set $L_{vj}\gets \mathrm G$ for all $v\in\mathcal V_{\mathrm{DISC}}$ and $1\leq j\leq N$
\State $E\gets n_{\mathrm{DISC}}N$ \Comment{events in one sweep}
\For{$t=1$ to $TE$}
    \State draw distinct $D,R$, and independently draw $V\in\mathcal V_{\mathrm{DISC}}$
    \If{$D\in\mathcal V_{\mathrm{ALLC}}$}
        \State $A^{\mathrm{int}}\gets \mathrm C$
    \ElsIf{$D\in\mathcal V_{\mathrm{ALLD}}$}
        \State $A^{\mathrm{int}}\gets \mathrm D$
    \Else
        \State draw $R_{DR}\sim\operatorname{Bernoulli}[\sigma(L_{DR})]$
        \State $A^{\mathrm{int}}\gets \mathrm C$ if $R_{DR}=1$; otherwise $A^{\mathrm{int}}\gets \mathrm D$
    \EndIf
    \State $A\gets A^{\mathrm{int}}$; if $A^{\mathrm{int}}=\mathrm C$, set $A\gets\mathrm D$ with probability $e_1$
    \State draw $R_{VR}\sim\operatorname{Bernoulli}[\sigma(L_{VR})]$
    \State draw prescribed assessment $Z$ from the norm entry $d_{A,R_{VR}}$
    \State reverse $Z$ with probability $e_2$
    \State $L_{VD}\gets T^+(L_{VD})$ if $Z$ is positive; otherwise $L_{VD}\gets T^-(L_{VD})$
    \If{$t\ge T_0E$ and $t-T_0E$ is divisible by $\Delta T E$}
        \State record the instantaneous statistics in Eqs.~\eqref{eq:si-finite-goodness}--\eqref{eq:si-finite-target-variance}
    \EndIf
\EndFor
\end{algorithmic}
\end{algorithm}

\subsection{Time unit, sampling and replication}

One event updates exactly one retained observer--target record. We define one Monte Carlo sweep as $E=n_{\mathrm{DISC}}N$ events, so that the mean number of update events per retained record is one per sweep. The default analyses run for $T=1{,}000$ sweeps. We discard the first $T_0=800$ sweeps and sample the remaining trajectory every $\Delta T=5$ sweeps. Each norm and parameter combination in the repeated depth analysis is simulated in 30 independent replicates. The base seed is 20260918; condition-specific seeds are generated by NumPy SeedSequence from the base seed, norm index (SC, SH, SJ, SS), $k_1$, $k_2$ and replicate index, and are saved with the output. Sampling includes both sweep 800 and sweep 1000, giving 41 post-burn-in measurements per run. A replicate-level estimate is the average of its post-burn-in samples; replicates, rather than successive samples within a replicate, are the units used for uncertainty intervals. The two extended SJ horizons are documented in Supplementary Note~\ref{sec:si-selection-robustness}.

All conditions use a fixed base seed and the same predetermined list of replicate seeds. Each condition is initialized and simulated independently. The saved configuration file reports every population, parameter, timing and seed value needed to reproduce an analysis.

\subsection{Mean reputation, disagreement and reputational separation}
\label{sec:si-finite-estimators}

At a sampled time, define
\[
s_{vj}=\sigma(L_{vj}),
\]
the conditional probability that observer $v$ reads target $j$ as good. Let $\mathcal O_j=\mathcal V_{\mathrm{DISC}}$ for every target, including its self-image if it is a discriminator, and write $m_j=|\mathcal O_j|=n_{\mathrm{DISC}}$. The expressed goodness of target $j$ is
\[
g_j=\frac{1}{m_j}\sum_{v\in\mathcal O_j}s_{vj}.
\]
For target strategy $I$, the mean probability of a good reputation is
\begin{equation}
g_I=\frac{1}{n_I}\sum_{j\in\mathcal V_I}g_j.
\label{eq:si-finite-goodness}
\end{equation}
This estimator averages the probabilities of the binary reputations that are actually consulted; it does not assign an arbitrary binary label to an intermediate internal state.

For a given target $j$, the probability that two distinct observers return different binary reputations in independent consultations is
\begin{equation}
d_j=
\binom{m_j}{2}^{-1}
\sum_{\substack{v,w\in\mathcal O_j\\v<w}}
\left[s_{vj}(1-s_{wj})+(1-s_{vj})s_{wj}\right].
\label{eq:si-finite-disagreement-target}
\end{equation}
The implementation evaluates this exact finite-observer expression as
\[
d_j=\frac{(m_j-1)S_j-S_j^2+Q_j}{\binom{m_j}{2}},\qquad
S_j=\sum_{v\in\mathcal O_j}s_{vj},\quad Q_j=\sum_{v\in\mathcal O_j}s_{vj}^2.
\]
For deterministic readouts with good fraction $x_j$, it reduces to $2m_jx_j(1-x_j)/(m_j-1)$, rather than the infinite-observer expression $2x_j(1-x_j)$. The original estimator already includes this correction; aligning the observer-pool description does not change its numerical implementation.
We define strategy-specific disagreement as
\begin{equation}
d_I=\frac{1}{n_I}\sum_{j\in\mathcal V_I}d_j.
\label{eq:si-finite-disagreement}
\end{equation}
Equations~\eqref{eq:si-finite-disagreement-target} and \eqref{eq:si-finite-disagreement} integrate over readout randomness exactly instead of adding Monte Carlo noise by drawing a display matrix. At the four endpoints $\alpha,\beta\in\{0,1\}$, every internal state has a deterministic binary readout. If $a_j$ of the $m_j$ observers then regard target $j$ as good, $d_j=2a_j(m_j-a_j)/[m_j(m_j-1)]$. A value $d_I=0$ denotes complete agreement; large values indicate that observers frequently express opposing reputations of the same targets. Mean goodness and disagreement are complementary: identical values of $g_I$ can arise from consensus within each target or from strong observer disagreement, so the reputation matrix cannot be summarized by $g_I$ alone.

To separate disagreement about a given target from differences among targets of the same strategy, target-to-target heterogeneity can additionally be quantified by
\begin{equation}
v_I=\frac{1}{n_I}\sum_{j\in\mathcal V_I}(g_j-g_I)^2.
\label{eq:si-finite-target-variance}
\end{equation}
Here $d_I$ measures observer disagreement conditional on the target, whereas $v_I$ measures target-to-target heterogeneity after observers have been averaged. The latter optional diagnostic is not used in the current selection rule or reported depth-grid outputs.

Finally, we quantify how well $\mathrm{DISC}$ observers separate cooperative or conditional behaviour from unconditional defection using
\begin{equation}
\Delta_{\mathrm{DISC}}=g_{\mathrm{DISC}}-g_{\mathrm{ALLD}},
\qquad
\Delta_{\mathrm{ALLC}}=g_{\mathrm{ALLC}}-g_{\mathrm{ALLD}}.
\label{eq:si-finite-separation}
\end{equation}
High overall goodness is not sufficient evidence of successful reputation formation, because it may also mean that unconditional defectors are treated as good. We therefore interpret an inertia condition jointly through $g_{\mathrm{DISC}}$, $g_{\mathrm{ALLC}}$, $g_{\mathrm{ALLD}}$, the two separation measures and $d_I$.

\subsection{Reputation matrices}

The matrix panels display the expected binary readout $s_{vj}$ for every retained $\mathrm{DISC}$ observer--target record at the end of one representative replicate. Rows are $\mathrm{DISC}$ observers and columns are targets, grouped as $\mathrm{DISC}$, $\mathrm{ALLC}$ and $\mathrm{ALLD}$. Self-records are included. The endpoint matrix panels are regenerated illustrative terminal configurations from a designated independent seed for each readout condition, not a replicate selected by its outcome. Their seeds and latent matrices are saved separately from the repeated depth-grid estimates. The matrix is a snapshot used to show the organization of private opinions; quantitative depth comparisons use the time-averaged replicate estimates defined above.

\subsection{Increasing the depth of reputation inertia}
\label{sec:si-depth-scan}

The one-step mechanism has one warning state and one suspicious state ($k_1=k_2=1$). To test how the amount and allocation of evidence affect recognition, we vary $k_1$ and $k_2$ independently over $0,\ldots,6$, with $\alpha=\beta=1$. The grid includes the binary baseline $(0,0)$ and full one-step inertia $(1,1)$. All population, error, timing and sampling parameters are held fixed.

All depth analyses use the direct-reset rule in Supplementary Fig.~\ref{fig:si-state-transitions}. A transition from established good to bad then requires $k_1+1$ consecutive negative assessments: any intervening positive assessment returns a warning state directly to $\mathrm G$. Similarly, a transition from established bad to good requires $k_2+1$ consecutive positive assessments, because any intervening negative assessment returns a suspicious state directly to $\mathrm B$. Increasing depth therefore changes the evidence threshold for overturning an established reputation. It is not merely a longer delay measured in events.

For every depth we report $g_I$, $d_I$, reputational separation, mean disagreement, and occupancy of the warning and suspicious states. The recognition--disagreement frontier and threshold-sensitivity analysis assess the robustness of the representative selections; their definitions are given below.

\subsection{Statistical summaries}

For each metric, the reported point estimate is the arithmetic mean of the 30 independent replicate means. We obtain percentile 95\% confidence intervals by resampling these replicate means with replacement 2,000 times and taking the 2.5th and 97.5th percentiles of the bootstrap distribution. No event-level or within-run samples are treated as independent replicates. The depth-analysis tables report the full two-dimensional grid, operational threshold sensitivity, and non-dominated sets. Selection uses the shortest qualifying architecture, not the largest separation irrespective of depth.

\subsection{Operational selection and robustness across architectures}
\label{sec:si-selection-robustness}

We use reputational separation $\Delta g=(g_{\mathrm{DISC}}+g_{\mathrm{ALLC}})/2-g_{\mathrm{ALLD}}$ and target-weighted mean disagreement $\bar d=\sum_I(n_I/N)d_I$. The working cutoffs $\Delta g\geq0.8$ and $\bar d\leq0.1$ are operational criteria used to identify compact architectures with high behavioural separation and low observer disagreement. They are not universal theoretical thresholds. Within the surveyed grid $0\leq k_1,k_2\leq6$, we first minimize $K=k_1+k_2$ among qualifying architectures, then break ties by greater separation and lower disagreement. This identifies a representative shortest qualifying architecture rather than a unique global optimum.

Selection robustness is evaluated over separation cutoffs from $0.70$ to $0.90$ in increments of $0.05$ and disagreement cutoffs from $0.050$ to $0.150$ in increments of $0.025$. We record the selected pair and total depth at each combination. We also report a conservative, pointwise-interval screen requiring the lower 95\% confidence bound on separation and upper bound on disagreement to satisfy the cutoffs. These intervals are not simultaneous confidence guarantees over the entire grid.

Independently of cutoffs, we identify architectures not dominated in the separation--disagreement plane, both over the surveyed grid and within each fixed $K$. Because an unconstrained two-outcome frontier does not penalize memory depth, we additionally report non-dominance in $(K,-\Delta g,\bar d)$. Frontier membership is descriptive and depends on the surveyed range and simulation uncertainty; it is not used to manufacture a uniquely preferred pair.

Independent runs, not within-run samples, are the uncertainty units. For a stationarity diagnostic we compare the early and late halves of each post-burn-in window and bootstrap their paired difference across runs. An interval wholly within $[-0.01,0.01]$ is recorded as a small-drift diagnostic for each of $\Delta g$ and $\bar d$. This practical diagnostic is not a proof of mixing or an additional recognition threshold; cases that fail require a longer observation window. Failure to detect improvement is not evidence of saturation: increasing repetitions can resolve progressively smaller effects. A formal saturation claim would require a prespecified practical effect size, a specified direction of depth increase and an equivalence analysis.

The representative pairs connect the main figures. General conclusions about evidence allocation should additionally be checked across neighboring qualifying architectures and the non-dominated sets at fixed total depth. Thus neither a single cutoff nor one selected pair is asked to establish generality.


\begin{table}[htbp]
\centering\small
\caption{Repeated estimates for the representative shortest qualifying architectures. Brackets show pointwise 95\% bootstrap confidence intervals across 30 independent replicate means.}
\label{tab:si-revalidated-architectures}
\begin{tabular}{llll}\toprule
Norm & $(k_1,k_2)$ & Separation & Disagreement\\\midrule
SC & $(2,0)$ & 0.9195 [0.9192, 0.9198] & 0.0644 [0.0641, 0.0646] \\
SH & $(3,0)$ & 0.9059 [0.9056, 0.9062] & 0.0694 [0.0691, 0.0696] \\
SJ & $(4,3)$ & 0.8806 [0.8801, 0.8812] & 0.0812 [0.0809, 0.0815] \\
SS & $(1,2)$ & 0.9174 [0.9171, 0.9177] & 0.0562 [0.0560, 0.0564] \\
\bottomrule\end{tabular}
\end{table}

At the default cutoffs, all four original pairs were retained in 1,000 of 1,000 bootstrap resamplings of the replicate-level data. This is conditional on the working thresholds and surveyed grid, not a probability of universal optimality. Across the 25 threshold combinations, the original pairs were selected in 12, 12, 6 and 14 cases for SC, SH, SJ and SS, respectively. Threshold changes therefore alter the representative pair. All four original pairs are non-dominated when total depth is included as a cost, whereas none lies on the unconstrained two-outcome separation--disagreement frontier. Very low disagreement can also arise from consensus on bad reputations; a two-outcome frontier alone does not identify a useful unique representative.

The standard 1,000-sweep horizon required extension at two SJ grid points. The $(4,5)$ condition was evaluated for 3,000 sweeps with 2,400 discarded, and $(3,5)$ for 50,000 sweeps with 40,000 discarded, retaining the five-sweep sampling interval and 30 independent runs. Both then met the small-drift diagnostic. These exceptions are recorded per output row. They do not change the default representative selections. Intermediate shorter-horizon data are retained for traceability.

\begin{table}[htbp]
\centering
\caption{\textbf{Default parameters for the finite-population reputation simulations.}}
\label{tab:si-finite-simulation-parameters}
\begin{tabular}{lll}
\toprule
Quantity & Symbol or option & Default value\\
\midrule
Population size & $N$ & 150\\
Strategy counts & $(n_{\mathrm{DISC}},n_{\mathrm{ALLC}},n_{\mathrm{ALLD}})$ & $(50,50,50)$\\
Implementation error & $e_1$ & 0.02\\
Assessment error & $e_2$ & 0.02\\
Initial internal state & $L_{vj}(0)$ & $\mathrm G$\\
Total duration & $T$ & 1,000 sweeps\\
Burn-in & $T_0$ & 800 sweeps\\
Sampling interval & $\Delta T$ & 5 sweeps\\
Independent replicates & --- & 30\\
Bootstrap resamples & --- & 2,000\\
Default depth range & --- & $0,\ldots,6$\\
Transition rule & --- & direct reset\\
\bottomrule
\end{tabular}
\end{table}




\begin{thebibliography}{1}
\bibitem{si-michelmata2024}
Michel-Mata, S., Kawakatsu, M., Sartini, J., Kessinger, T. A., Plotkin, J. B. \& Tarnita, C. E. The evolution of private reputations in information-abundant landscapes. \textit{Nature} \textbf{634}, 883--889 (2024). doi:10.1038/s41586-024-07977-x.
\bibitem{si-schmid2023}
Schmid, L., Ekbatani, F., Hilbe, C. \& Chatterjee, K. Quantitative assessment can stabilize indirect reciprocity under imperfect information. \textit{Nature Communications} \textbf{14}, 2086 (2023). doi:10.1038/s41467-023-37817-x.
\bibitem{si-nowak2005}
Nowak, M. A. \& Sigmund, K. Evolution of indirect reciprocity. \textit{Nature} \textbf{437}, 1291--1298 (2005). doi:10.1038/nature04131.
\bibitem{si-hilbe2018}
Hilbe, C., Schmid, L., Tkadlec, J., Chatterjee, K. \& Nowak, M. A. Indirect reciprocity with private, noisy, and incomplete information. \textit{Proceedings of the National Academy of Sciences USA} \textbf{115}, 12241--12246 (2018). doi:10.1073/pnas.1810565115.
\end{thebibliography}



\end{document}
